\documentclass[11pt]{article}

\usepackage[a4paper,margin=27mm]{geometry}
\usepackage{fontspec}
\setmainfont{Latin Modern Roman}
\setsansfont{Latin Modern Sans}
\setmonofont{DejaVu Sans Mono}[Scale=MatchLowercase]
\usepackage{microtype}
\usepackage{amsmath,amssymb,amsthm,mathtools}
\usepackage{booktabs}
\usepackage{array}
\usepackage{longtable}
\usepackage{float}
\usepackage{enumitem}
\usepackage{xcolor}
\usepackage[hidelinks]{hyperref}
\def\UrlBreaks{\do\-\do\/\do\.}
\setlength{\emergencystretch}{2em}
\hbadness=10000
\hfuzz=3pt

\newtheoremstyle{upright}
  {6pt}{6pt}{\normalfont}{}{\bfseries}{.}{.5em}{}
\theoremstyle{upright}
\newtheorem{theorem}{Theorem}[section]
\newtheorem{proposition}[theorem]{Proposition}
\newtheorem{corollary}[theorem]{Corollary}
\newtheorem{lemma}[theorem]{Lemma}
\theoremstyle{definition}
\newtheorem{definition}[theorem]{Definition}
\theoremstyle{remark}
\newtheorem{remark}[theorem]{Remark}

\newcommand{\F}{\mathbb F}
\newcommand{\PP}{\mathbb P}
\newcommand{\PRS}{\operatorname{PRS}}
\newcommand{\PGL}{\operatorname{PGL}}
\newcommand{\PGammaL}{\mathrm{P}\Gamma\mathrm{L}}
\newcommand{\Sym}{\operatorname{Sym}}
\newcommand{\Span}{\operatorname{span}}
\newcommand{\Tr}{\operatorname{Tr}}
\newcommand{\Res}{\operatorname{Res}}
\newcommand{\Hess}{\operatorname{Hess}}
\newcommand{\cC}{\mathcal C}
\newcommand{\cM}{\mathcal M}
\newcommand{\cN}{\mathcal N}
\newcommand{\cP}{\mathcal P}
\newcommand{\cU}{\mathcal U}
\newcommand{\cW}{\mathcal W}
\newcommand{\NRC}{\mathcal R}

\title{Projective Reed--Solomon syndromes beyond redundancy four:
deep holes, coherent polar flags, and modular carriers}
\author{Tavis Rudd}
\date{July 2026}

\begin{document}
\maketitle

\begin{abstract}
A projective Reed--Solomon syndrome determines a Hankel linear system
of binary forms.  We call the direction split-free when that system
contains no completely split squarefree member.  This reformulation
gives a complete classification of the redundancy-five deep-hole
directions over every prime power $q\geq7$; the generic case is
controlled by an integral off-diagonal fiber-square curve of arithmetic
genus one.

For higher redundancy we introduce coherent polar flags: iterated
contractions in which every selected factor remains a marked forbidden
root.  After separating the level-specific contained-component problem
from the transverse point-count argument, this method gives complete
all-field classifications at redundancies six and seven; the
redundancy-seven split-free list is a deep-hole classification whenever
the covering radius is six, in particular for $q\geq11$.  The analogous transverse
arguments at redundancies eight and nine leave only the persistent
families for $q\geq43$ and $q\geq53$, conditional on their printed
component and slice gates.

In characteristic two we determine the geometric degeneracy strata of
an ordered Hessian pullback and isolate the remaining root-compatible
component and base-selection gates.  We also analyze the
power-of-two Lucas arithmetic and the distinguished degree-nine
$e_7$ orbit.  The remaining degree-nine carrier strata and
redundancy-ten deep holes are not classified.
\end{abstract}

\section{Introduction}

Let $\PRS(k)$ denote the projective Reed--Solomon code of length $q+1$
obtained by evaluating homogeneous binary forms of degree $k-1$ on
$\PP^1(\F_q)$.  If the redundancy is
\[
  r=q+1-k,
\]
a parity-check matrix has as columns the $\F_q$-points of the
degree-$(r-1)$ normal rational curve
\[
  \NRC_{r-1}=\{(1,t,\ldots,t^{r-1}):t\in\F_q\}
              \cup\{(0,\ldots,0,1)\}\subset\PP^{r-1}.
\]
Call a projective syndrome direction \emph{split-free} if it is not
contained in the span of $r-2$ distinct columns.  When
$\rho(\PRS(q+1-r))=r-1$, the split-free directions are precisely the
projective syndrome directions of deep-hole cosets.  The correspondence between deep holes,
one-symbol MDS extensions, and uncovered points of a normal rational
curve was developed in \cite{Kaipa2017,ZWK2020}.  Zhang, Wan, and Kaipa
classified projective Reed--Solomon deep-hole directions through
redundancy four and identified redundancy five as the next case
\cite{ZWK2020}.  Our first theorem gives the corresponding all-field
redundancy-five classification.

\begin{table}[ht]
\centering
\small
\caption{Notation used throughout the paper.}
\label{tab:notation}
\begin{tabular}{cl}
\toprule
Symbol & Meaning\\
\midrule
$q$ & field size; all numerical thresholds are thresholds on prime powers\\
$k$ & code dimension in $\PRS(k)$\\
$r=q+1-k$ & redundancy\\
$n=r-1$ & divided-power degree of a syndrome representative\\
$E$ & a two-dimensional vector space over the current base field\\
$f\in\Gamma^nE$ & divided-power representative of a projective syndrome\\
$W_f\subset\Sym^{n-1}E^\vee$ & Hankel/apolar witness system annihilated by $f$\\
$g,\delta,\kappa$ & genus, deletion degree, and point-bound correction\\
$b,c$ & transverse bad-carrier and marker-collision degrees\\
$d=2^s$ & power-of-two Lucas degree; used only in Section~\ref{sec:lucas}\\
\bottomrule
\end{tabular}
\end{table}

\begin{definition}[Working terminology]\label{def:terminology}
A \emph{carrier} is a closed subvariety of a syndrome space on which
the generic transverse splitting argument degenerates.  A
\emph{persistent carrier} is a catalecticant component propagated by
ordinary contraction; a \emph{modular} or \emph{Lucas carrier} is a
component created by a positive-characteristic contraction kernel.
A direction is \emph{shallow} when it has a split squarefree witness.
A pullback is \emph{root-compatible} when the chosen contraction
factor remains a marked forbidden root on every lower incidence.  An
\emph{ordered-root cover} remembers an ordering of the roots over the
squarefree splitting locus.  A field is \emph{code-admissible} for a
displayed statement when it satisfies both that statement's field
bound and its covering-radius hypothesis.
\end{definition}

\paragraph{Status of this draft.}
This integrated version freezes the theorem spine and the computational
trust boundaries.  It is a research-announcement draft: the
degree-specific proof expansions and immutable external supplement
listed in the proof ledger remain release gates.  In particular, this
version is not described as a proof-complete submission.

\begin{theorem}[Complete redundancy-five classification]\label{thm:r5}
Let $q\geq7$ be a prime power.  The covering radius of $\PRS(q-4)$ is
$4$.  Its projective deep-hole syndromes are exactly the following
disjoint families:
\begin{enumerate}[label=\textup{(\roman*)}]
\item the tangent family $\mathrm T$, of size $q(q+1)$;
\item the conjugate-secant family $\mathrm S$, of size
      $q(q^2-1)/2$;
\item if $\operatorname{char}\F_q\ne3$, the rational osculating-pair
      family $\mathrm O^+$ when $q\equiv2\pmod3$, of size $q(q+1)/2$,
      or the conjugate osculating-pair family $\mathrm O^-$ when
      $q\equiv1\pmod3$, of size $q(q-1)/2$;
\item in characteristic three, one $\PGL_2(q)$-fixed nucleus point and
      one wild Artin--Schreier orbit $\mathrm W$ of size $(q^2-1)/2$;
\item the sporadic $\PGL_2(q)$-orbits in Table~\ref{tab:r5sporadic},
      occurring exactly for
      $q\in\{7,8,9,11,13,17,19\}$.
\end{enumerate}
There are no other deep-hole syndromes.  The three corresponding
nonsporadic totals are
\[
\frac{q^2(q+3)}2,\qquad
\frac{q(q+1)(q+2)}2,\qquad
\frac{q^3+3q^2+q+1}{2},
\]
according as $q\equiv1\pmod3$ with characteristic not three,
$q\equiv2\pmod3$, or $\operatorname{char}\F_q=3$.
\end{theorem}

The main mechanism is a characteristic-free Hankel dictionary.  A
syndrome at redundancy $r$ annihilates a linear system of binary forms
of degree $r-2$.  It is split-free if and only if this system contains
no completely split squarefree form; it is a deep-hole direction when,
in addition, $\rho(\PRS(q+1-r))=r-1$.  At redundancy five the system is a
pencil of cubics.  Its induced degree-three map of projective lines is
either cyclic or has geometric monodromy $S_3$.  Cyclic descent gives
the arithmetic families in Theorem~\ref{thm:r5}; in the $S_3$ case the
off-diagonal fiber square is an integral curve of arithmetic genus one,
and the singular-curve Weil bound \cite{AubryPerret1995} forces a split
fiber for $q\geq23$.  The certified finite computations in
Section~\ref{sec:verification} handle the remaining fields.

Starting at redundancy six, the Hankel system is no longer a pencil.
The useful replacement is the coherently parameterized line of
contractions obtained by choosing one prospective root.  Iteration
produces a \emph{polar flag}; every chosen root remains marked because
a lower witness containing it would lift with multiplicity two.  The
pointed coherence is essential: at $q=19$ an individual contracted
fiber has no split witness avoiding its marker, but no consecutive-row
polar line is contained in its orbit.

The resulting structural theorem has three outputs.

\begin{theorem}[Classification results and explicit gates beyond five]
\label{thm:spine}
The following statements hold.
\begin{enumerate}[label=\textup{(\alph*)}]
\item For every prime power $q\geq7$, the projective deep-hole
      directions of $\PRS(q-5)$ are exactly the persistent tangent and
      conjugate-secant families, the exceptional orbits listed for
      $q=7,8,9,11,13$, and, when $q=2^m$ with odd $m\geq5$, one
      characteristic-two nucleus orbit.
\item The displayed finite domains through $q=32$ have the certified
      split-free classification stated in Section~\ref{sec:r67}.
      The same list continues over every field: for $q\geq13$ it is the
      persistent families plus one central nucleus point when
      $q=2^m$ and $m$ is odd.  The list is a deep-hole classification
      only when $\rho(\PRS(q-6))=6$, in particular for $q\geq11$.
\item Conditional on the degree-seven contained-component gate and the
      degree-eight six-slice/base-selection gates printed in
      Section~\ref{sec:fixed}, the projective deep-hole directions of
      $\PRS(q-7)$ for $q\geq43$ and of $\PRS(q-8)$ for $q\geq53$ are
      exactly the persistent tangent and conjugate-secant families.
      In either case their total size is $q(q+1)^2/2$.
\end{enumerate}
At redundancy $r$, the semilinear orbits in the persistent
conjugate-secant family are parameterized by
\[
       T/T^{\,r-1}\quad\text{modulo inversion and Frobenius},
\]
where
$T=\{z\in\F_{q^2}^{\times}:z^{q+1}=1\}$ is the norm-one torus and
Frobenius is coefficient Frobenius.  On the tangent family the
unipotent coordinate transforms as $z\mapsto z+(r-1)u$; consequently
its orbit decomposition changes precisely when
$\operatorname{char}\F_q$ divides $r-1$.
\end{theorem}

The three parts have different scopes.  Parts~(a) and~(b) are
unconditional all-field syndrome classifications, but part~(b)
becomes a deep-hole classification only when
$\rho(\PRS(q-6))=6$.  Part~(c) isolates the remaining mathematical
gates rather than treating them as certificate inputs and makes no
classification claim below its displayed high-field thresholds.

\paragraph{Relation to previous work.}
The tangent and conjugate-secant families and the coding/normal-rational
curve dictionary come from \cite{Kaipa2017,ZWK2020}; the covering-radius
input in the high-rate range is due to Seroussi and Roth
\cite{SeroussiRoth1986}.  The broader arc/MDS and code-automorphism
background is surveyed or established in
\cite{BallLavrauw2019,Dur1987}.  Classical apolarity, catalecticants,
and binary-form rank strata are treated in
\cite{IarrobinoKanev1999,ComasSeiguer2011}; our use of divided powers
keeps the small-characteristic contraction formulas explicit.
Binary-quartic orbit theory in characteristic different from two and
three \cite{KPP2025}, factorization statistics in linear families
\cite{CMP2017}, normal-rational-curve nuclei
\cite{GmainerHavlicek2013}, and the general Frobenius/monodromy
semantics of splitting families \cite{Wang2026} are inputs.  The first
two references identify deep syndromes with uncovered points of a
normal rational curve and settle the lower-redundancy families; they do
not classify the redundancy-five cubic pencils.  The orbit-theoretic
and factorization references supply normal forms and distribution
tools, but not the coherent marked contraction used here.  The nucleus
results locate characteristic-dependent kernels, while the present
work computes which of their coherent lifts meet the PRS splitting
incidence.  Standard function-field, Artin--Schreier, and linearized
polynomial background is taken from
\cite{Stichtenoth2009,LidlNiederreiter1996}; the characteristic-two
orbit and component exhaustions themselves are proved internally, not
imported from the quartic-orbit reference.  Thus the new assertions are separated into the
redundancy-five through seven classifications, the conditional
transverse induction mechanism, the fixed-level high-field
applications, and the characteristic-two contained-carrier results.
Claim-by-claim novelty and dependency boundaries are recorded in the
electronic ledger; no database-search procedure is part of the
mathematical argument.

\paragraph{Organization.}
Section~\ref{sec:overview} first gives a proof-free guide to the
argument, and Section~\ref{sec:dictionary} gives the Hankel dictionary.
Section~\ref{sec:r5} proves Theorem~\ref{thm:r5}.
Sections~\ref{sec:polar}--\ref{sec:fixed} develop coherent polar
induction and prove Theorem~\ref{thm:spine}.
Section~\ref{sec:hessian} gives the characteristic-two ordered-Hessian
replacement.  Section~\ref{sec:lucas} treats Lucas-carrier arithmetic.
Section~\ref{sec:verification} states the computational and formal trust
boundary, and Section~\ref{sec:boundary} records the exact open boundary.

\section{Reading map}\label{sec:overview}

Here is the argument in one pass.  A syndrome \(f\) determines a linear
system \(W_f\) of binary forms through its two-row Hankel matrix.  The
Hankel criterion says that a completely split squarefree form in \(W_f\)
is a witness that \(f\) lies in the span of the required number of distinct
normal-rational-curve points.  Thus the geometric problem is to find such a
form; only after the separate covering-radius gate does its absence mean that
\(f\) is a deep-hole syndrome.

At redundancy five, \(W_f\) is a pencil of cubics.  Its off-diagonal fibre
square records ordered pairs of distinct roots lying in one pencil member.
The exact count in Proposition~\ref{prop:r5-count} turns rational
points on that genus-one curve into split cubics and makes R5 the terminal
model for every later level.

At higher redundancy, choose a root \(\lambda\) as a marker and contract the
syndrome.  A split witness in the lower system lifts after multiplication by
\(\lambda\), and it stays squarefree precisely when it avoids that marker.
As \(\lambda\) varies, the contracted syndromes trace a polar line.  There
are then two branches.  If the line is transverse to the lower bad locus,
only finitely many markers and points on the lower splitting cover must be
discarded; a rational-point bound supplies a witness.  If the line is
contained, the component calculation places the original syndrome in the
persistent catalecticant locus or a characteristic-dependent Lucas carrier.
Figure~\ref{fig:polar-flow} displays this single recursion.

\begin{figure}[H]
\centering
\[
\begin{array}{rcl}
\text{coding}
 &:& [f]\longmapsto W_f\supset g_{\mathrm{split}}
      \longmapsto [f]\text{ is shallow},\\[4pt]
\text{recursion}
 &:& [f]\xrightarrow{\ \text{contract at }\lambda\ }
      \text{lower splitting problem},\\[2pt]
&&\begin{array}{rcl}
\text{contained}&\longmapsto&\text{persistent or Lucas carrier},\\
\text{transverse}&\longmapsto&
 \text{lower witness }g,\ \lambda\nmid g
 \longmapsto \lambda g\in W_f .
\end{array}
\end{array}
\]
\caption{The conceptual route through the paper.  The first line changes
from coding language to a splitting problem.  Marked contraction then
separates the contained branch, classified by the carrier theorem, from the
transverse branch, where finite avoidance and point counting produce the
lower witness \(g\).}
\label{fig:polar-flow}
\end{figure}

The proof has one recursive spine and six evaluated levels.
Table~\ref{tab:dependency-map} shows what changes from one level to the next;
its deletion budget is the number of points excluded from the terminal cover.
Table~\ref{tab:classification-status} separately records which geometric
classifications promote to code deep-hole classifications.

\begin{table}[H]
\centering
\small
\caption{Results and mechanisms.  The threshold column records the first
field range in which the stated geometric argument closes; bounded
certificates handle only the listed lower fields.}
\label{tab:dependency-map}
\setlength{\tabcolsep}{3pt}
\begin{tabular}{p{0.07\textwidth}p{0.24\textwidth}p{0.40\textwidth}
                p{0.22\textwidth}}
\toprule
Level & New obstruction & Main geometric mechanism & Closing gate\\
\midrule
R5 & cubic splitting & cubic pencil and off-diagonal \((2,2)\) fibre square;
exact split-witness count & branch correction at most \(12\), or \(6\) in
characteristic two; geometric threshold \(20\) and seven-field certificate\\
R6 & one retained marker & one contraction to the same R5 curve & deletion
budget \(19\); geometric threshold \(29\) and finite bridge\\
R7 & two retained markers & two coherent contractions to the same R5 curve &
deletion budget \(25\); geometric threshold \(37\) and finite bridge\\
R8 & three retained markers & the same integral \((2,2)\) terminal cover &
deletion budget \(30\); \(q\geq43\)\\
R9 & four markers and a characteristic-seven carrier & the terminal cover
off the carrier; a residual-quadratic genus-one slice on it & deletion budget
\(36\); \(q\geq53\), with the first relevant characteristic-seven field
handled geometrically\\
R10 & five markers and the first higher Lucas carrier & the terminal cover
in odd characteristic, an ordered-Hessian slice in characteristic two,
and a final-pair Artin--Schreier criterion on the carrier & deletion
budget \(42\); odd \(q\geq59\), even \(q\geq64\), plus the stated binary
carrier certificates\\
\bottomrule
\end{tabular}
\end{table}

\begin{table}[H]
\centering
\small
\caption{Classification status.  A split-free equality becomes a code
deep-hole equality only after the radius gate.}
\label{tab:classification-status}
\setlength{\tabcolsep}{3pt}
\begin{tabular}{p{0.07\textwidth}p{0.15\textwidth}
                p{0.29\textwidth}p{0.29\textwidth}
                p{0.12\textwidth}}
\toprule
Red. & Field range & Split-free geometry & Deep holes and radius
& Gap\\
\midrule
\(5\) & \(q\geq7\) & Complete & Complete; \(\rho=4\) & None\\
\(6\) & \(q\geq7\) & Complete & Complete; \(\rho=5\) & None\\
\(7\) & \(q\geq7\) & Complete & Complete for \(q\geq11\) & Radius at
\(7,8,9\)\\
\(8\) & \(q\geq43\) & Persistent-only & Complete; \(\rho=7\) & Below range\\
\(9\) & \(q\geq53\) & Persistent-only & Complete; \(\rho=8\) & Below range\\
\(10\) & \(q\geq59\) & Persistent-only & Complete; \(\rho=9\) & Below range\\
\(r\geq6\) & \(q\geq Q_r\) & Contained in persistent plus maximal Lucas;
persistent-only if \(p>r-1\), conditional on the intermediate lower
packages & Complete under that hypothesis if \(p>r-1\) & Stagewise packages;
Lucas arithmetic in small characteristic\\
\bottomrule
\end{tabular}
\end{table}

Here \(Q_r=6r-15+\lfloor2\sqrt{6r-17}\rfloor\).  The R7 finite record
is exposed only as a four-field exception delta beyond the uniform
\(T/T^6\) spine: the full fourteen-field ledger remains in the supplement.


\section{The syndrome--Hankel dictionary}\label{sec:dictionary}

Let $E$ be a two-dimensional vector space.  We use divided powers for
syndromes and ordinary symmetric powers for root forms.  This avoids
invalid binomial rescalings in small characteristic.  For
$f=(a_0:\cdots:a_{r-1})\in\PP(\Gamma^{r-1}E)$, write
$e_0,\ldots,e_{r-1}$ for the divided-power coordinate basis and
$\nu(t)=(1,t,\ldots,t^{r-1})$, with
$\nu(\infty)=e_{r-1}$.  The perfect divided-power/symmetric-power
pairing is characterized by coefficient extraction.  Define
\[
 H_f=
 \begin{pmatrix}
 a_0&a_1&\cdots&a_{r-2}\\
 a_1&a_2&\cdots&a_{r-1}
 \end{pmatrix},\qquad
 W_f=\ker H_f\subset\Sym^{r-2}E^\vee.
\]

\begin{lemma}[Hankel criterion]\label{lem:hankel}
Let $\lambda_1,\ldots,\lambda_{r-2}$ be distinct points of
$\PP(E^\vee)(\F_q)$ and put $g=\prod_i\lambda_i$.  Then
\[
 f\in\Span\bigl(\nu(\lambda_1),\ldots,\nu(\lambda_{r-2})\bigr)
 \quad\Longleftrightarrow\quad g\in W_f.
\]
The same statement with repeated factors and Hasse derivatives gives
the corresponding osculating spans.  Both statements commute with
$\PGammaL_2(q)$.
\end{lemma}

\begin{proof}
On an affine chart, $g$ is the characteristic polynomial of an
order-$(r-2)$ recurrence.  The length-$r$ solutions are spanned by the
geometric sequences $\nu(\lambda_i)$; their independence is the
Vandermonde determinant.  A root at infinity gives the reversed-chart
recurrence.  For repeated roots, the confluent Vandermonde basis is
formed by Hasse derivatives, so the same argument works in every
characteristic.  Equivariance follows because substitution transports
both the normal rational curve and the recurrence.
\end{proof}

\begin{corollary}[Split-free criterion]\label{cor:splitfree}
Assume $\rho(\PRS(q+1-r))=r-1$.  A projective syndrome $f$ is deep if
and only if $W_f$ contains no completely split squarefree form of
degree $r-2$ over $\F_q$.
\end{corollary}

We call the latter property \emph{split-free} even when the
covering-radius hypothesis has not been established.  Thus every deep
syndrome in the radius-$r-1$ range is split-free, while a small-field
split-free census is not silently promoted to a code deep-hole
classification.

The persistent locus is already visible here.  If every member of
$W_f$ has a common quadratic factor, then that factor is either a
rational square, an irreducible quadratic, or a product of two distinct
rational factors.  These are respectively the tangent deep family, the
conjugate-secant deep family, and a shallow rational secant.  The common
factor cannot grow with redundancy: since $\dim W_f=r-3$,
\[
 r-3\leq r-1-\deg\gcd(W_f),
\]
so its degree is at most two.


\section{Redundancy five: exceptional cubic covers}\label{sec:r5}

Here $f=(a_0:\cdots:a_4)$ and $W_f$ is the pencil annihilated by
\[
 \begin{pmatrix}a_0&a_1&a_2&a_3\\a_1&a_2&a_3&a_4\end{pmatrix}.
\]
These are divided-power coordinates; the coefficients of a cubic in
$W_f$ are ordinary symmetric-power coordinates.  Factoring the formal
quartic $\sum a_iT^iU^{4-i}$ would therefore be noninvariant in
characteristics two and three.

\begin{proposition}[Radius gate]\label{prop:r5-radius}
For every prime power $q\geq7$,
\[
       \rho(\PRS(q-4))=4.
\]
\end{proposition}

\begin{proof}
The coding/arc criterion identifies radius four with completeness of
the quartic normal rational curve.  Seroussi--Roth prove completeness
when
\[
       k\geq\left\lfloor\frac{q-1}{2}\right\rfloor
\]
apart from the even-field dimensions $k=2,q-2$
\cite{SeroussiRoth1986}.  Here $k=q-4$, the inequality holds for
$q\geq7$, and neither exceptional dimension occurs.  Thus
Corollary~\ref{cor:splitfree} applies.  The finite certificate
independently checks the same conclusion only in its declared finite
domain; it is not the all-field radius proof.
\end{proof}

\subsection{Gcd strata}

\begin{proposition}[Quadratic gcd]\label{prop:r5-gcd2}
If $\deg\gcd(W_f)=2$, then
\[
       W_f=Q\langle T,U\rangle.
\]
If $Q$ is split squarefree, $f$ is shallow.  If $Q$ is irreducible,
$f$ belongs to the conjugate-secant family; if $Q=\lambda^2$, it
belongs to the tangent family.  Conversely every point in those
families has the displayed gcd.  Their deep-family sizes are
\[
       |\mathrm T|=q(q+1),\qquad
       |\mathrm S|=\frac{q(q^2-1)}2.
\]
\end{proposition}

\begin{proof}
The equality follows by dimension.  A split $Q$ and a linear cofactor
avoiding its two roots give a split squarefree cubic.  An irreducible
$Q$ survives in every member, while a square $Q$ forces a repeated
root in every member.  Lemma~\ref{lem:hankel}, including its confluent
and conjugate-root forms, identifies the corresponding syndromes with
the secant and tangent spans.  The sizes are the established
tangent/conjugate-secant counts from the lower-redundancy
classification \cite{Kaipa2017,ZWK2020}.
\end{proof}

\begin{proposition}[Linear gcd]\label{prop:r5-gcd1}
Suppose $\deg\gcd(W_f)=1$.  Then $f$ is shallow for $q\geq8$.  At
$q=7$ the same conclusion is part of Certificate R5.
\end{proposition}

\begin{proof}
Write $W_f=\lambda_0V$ with $V$ a basepoint-free quadratic pencil.  If
its degree-two map $\psi:\PP^1\to\PP^1$ is separable, its rational deck
involution gives $q+1$ ordered same-fiber pairs.  Delete at most two
fixed points, four points over branch values, and the at most two pairs
whose selected member contains $\lambda_0$.  For $q\geq8$ a surviving
pair gives a split quadratic whose roots avoid $\lambda_0$, and hence a
split squarefree cubic in $W_f$.  In characteristic two an inseparable
quadratic pencil is equivalent to $\langle T^2,U^2\rangle$; the two
overlapping Hankel equations then force $f$ onto the normal rational
curve, contrary to the rank-two pencil hypothesis.
\end{proof}

\subsection{The trivial-gcd cubic cover}

Choose a basis $c_1,c_2$ of a trivial-gcd pencil and put
\[
 X_f=\{([u:v],t):uc_1(t)+vc_2(t)=0\}
       \subset\PP^1\times\PP^1.
\]

\begin{proposition}[Cubic incidence and off-diagonal fiber square]
\label{prop:r5-incidence}
The curve $X_f$ is geometrically integral of bidegree $(1,3)$ and
arithmetic genus zero.  Projection to the root coordinate is
birational, while projection to the pencil coordinate gives a
degree-three morphism
\[
       \phi_f:\PP^1\longrightarrow\PP^1.
\]
If $\phi_f$ is separable, division of
\[
 c_1(t_1)c_2(t_2)-c_2(t_1)c_1(t_2)
\]
by the diagonal bracket defines a residual curve
$Y_f\subset\PP^1\times\PP^1$ of bidegree $(2,2)$ and arithmetic genus
one.  Its geometric components are the geometric-monodromy orbits on
ordered pairs of distinct roots.
\end{proposition}

\begin{proof}
A component of $X_f$ of bidegree $(0,e)$ would be a common factor of
$c_1,c_2$; hence none exists.  The bidegree genus formula gives genus
zero.  A root determines the unique pencil member containing it, so the
root projection is birational.  For a separable map the diagonal factor
in the fiber-square determinant has multiplicity one.  Removing its
bidegree $(1,1)$ from the fiber square leaves bidegree $(2,2)$, whose
arithmetic genus is one.  The standard fiber-square/monodromy
correspondence identifies its components with orbits on distinct
ordered sheets.
\end{proof}

\begin{lemma}[Cyclic stratum]\label{lem:cyclic}
Suppose $\phi_f$ is geometrically cyclic.
\begin{enumerate}[label=\textup{(\roman*)}]
\item In characteristic different from three there are two ramification
      points.  If they are rational, the corresponding osculating-plane
      intersection is deep exactly when $q\equiv2\pmod3$; if they are
      conjugate, it is deep exactly when $q\equiv1\pmod3$.
\item In characteristic three all osculating planes meet in
      $n=(0:0:1:0:0)$, whose pencil is
      $\langle T^3,U^3\rangle$; this fixed point is deep.
\item The remaining wild characteristic-three pencils have the form
      $\langle U^3,T^3+aTU^2\rangle$.  They are deep exactly when $-a$
      is nonsquare, forming one orbit of size $(q^2-1)/2$.
\end{enumerate}
\end{lemma}

\begin{proof}
Assume first that the characteristic is not three.  Riemann--Hurwitz
gives two totally ramified points $t_1,t_2$.  The confluent Hankel
criterion places $f$ in
$\operatorname{Osc}(t_1)\cap\operatorname{Osc}(t_2)$; after transporting
the pair to $\{0,\infty\}$, the pencil is
$\langle T^3,U^3\rangle$ and a geometric deck generator is
$t\mapsto\zeta t$.

For a rational ramification pair, Frobenius fixes the points and the
deck map descends exactly when $\zeta^q=\zeta$, namely
$q\equiv1\pmod3$.  Thus its complement is deep exactly for
$q\equiv2\pmod3$.  The setwise stabilizer of the ordered normal form is
the split dihedral group of order $2(q-1)$, so the orbit size is
\[
       \frac{q(q^2-1)}{2(q-1)}=\frac{q(q+1)}2.
\]
For a conjugate pair Frobenius interchanges the fixed points and
inverts the multiplier; descent is instead $q\equiv2\pmod3$.  The
nonsplit dihedral stabilizer has order $2(q+1)$ and the orbit size is
$q(q-1)/2$.

Now suppose the characteristic is three.  The second Hasse derivative
of
\[
       \nu(t)=(1,t,t^2,t^3,t^4)
\]
is $(0,0,1,3t,6t^2)=e_2$, including in the reversed infinity chart.
Thus every osculating plane contains $n=[e_2]$.  Substitution shows that
$n$ is $\PGL_2(q)$-fixed and
$W_n=\langle T^3,U^3\rangle$, so it is deep.

A remaining separable cyclic cover has one wild totally ramified
point.  Sending it to infinity and comparing coefficients under a
nontrivial translation gives
\[
 W_f=\langle U^3,T^3+\alpha T^2U+\beta TU^2\rangle,\qquad
 \alpha=0,\quad \kappa^2=-\beta.
\]
It is therefore represented by
\[
 f_a=e_2-ae_4,\qquad
 W_{f_a}=\langle U^3,T^3+aTU^2\rangle.
\]
The affine members are $z^3+az+s$.  Such a member has three rational
roots exactly when the additive map $z\mapsto z^3+az$ has a nonzero
rational kernel, equivalently when $-a$ is a square.  Hence the
nonsquare parameters form the deep orbit.  Its stabilizer is
$t\mapsto\pm t+\beta$, of order $2q$, and orbit--stabilizer gives
$(q^2-1)/2$.
\end{proof}

\begin{lemma}[$S_3$ stratum]\label{lem:s3}
If $\phi_f$ has geometric monodromy $S_3$ and $q\geq23$, then the pencil
contains a completely split squarefree cubic.
\end{lemma}

\begin{proof}
By Proposition~\ref{prop:r5-incidence}, components of $Y_f$ correspond
to monodromy orbits on distinct ordered roots.  The $S_3$ action is
two-transitive, so $Y_f$ is geometrically integral and has arithmetic
genus one.  Aubry--Perret gives
\[
       \#Y_f(\F_q)\geq q-2\sqrt q.
\]
The diagonal costs
\[
       Y_f\cdot\Delta=(2,2)\cdot(1,1)=4.
\]
Riemann--Hurwitz gives different degree four for the separable cubic
map, hence at most four branch values.  A singular cubic has at most
two distinct roots and at most two ordered off-diagonal pairs, so the
branch deletion costs at most eight.  If
$q-2\sqrt q>12$, a rational off-diagonal unramified point remains.  It
gives two distinct rational roots of a squarefree cubic, whose third
root is rational.  The inequality holds for every prime power
$q\geq23$.
\end{proof}

\begin{proposition}[Exhaustion and finite bridge]\label{prop:r5-bridge}
The preceding cases exhaust every redundancy-five pencil for
$q\geq23$.  Certificate R5 closes exactly
\[
       q\in\{7,8,9,11,13,16,17,19\}.
\]
It finds sporadic $S_3$-stratum orbits exactly at
$q=7,8,9,11,13,17,19$ and none at $q=16$.
\end{proposition}

\begin{proof}
The gcd degree is zero, one, or two.  Propositions
\ref{prop:r5-gcd2} and \ref{prop:r5-gcd1} settle the positive-gcd
cases.  In the trivial-gcd case the cubic map is inseparable or
separable.  The inseparable case occurs only in characteristic three
and the overlap equations give the all-cube nucleus pencil.  A
separable transitive cubic has geometric monodromy $C_3$ or $S_3$,
settled by Lemmas~\ref{lem:cyclic} and \ref{lem:s3}.  The stated finite
residue, canonical representatives, stabilizers, histograms, and
Frobenius links are exactly the domain of Certificate R5.
\end{proof}

\begin{table}[ht]
\centering
\caption{Complete sporadic orbit inventory at redundancy five.  Entries
are $\PGL_2(q)$ orbit sizes; exponents denote multiplicity.}
\label{tab:r5sporadic}
\begin{tabular}{ccl}
\toprule
$q$ & number of sporadic points & orbit sizes\\
\midrule
7  & 644 & $28,112,168^3$\\
8  & 756 & $252^3$\\
9  & 900 & $180,360^2$\\
11 & 990 & $330,660$\\
13 & 728 & $182,546$\\
17 & 1224& $1224$\\
19 & 570 & $570$\\
\bottomrule
\end{tabular}
\end{table}

The orbit sizes in Table~\ref{tab:r5sporadic} are a compact inventory,
not orbit identifiers.  The electronic C491 certificate records for
each row a canonical normalized representative, stabilizer, complete
cubic-pencil member histogram, and Frobenius image; equality with one
of those canonical $\PGL_2$ orbits is the finite classification
criterion.  Thus repeated sizes denote distinct certified orbits.

\section{Coherent polar induction}\label{sec:polar}

\subsection{Spaces, markers, and bad schemes}

Fix a field $k$, a two-dimensional $k$-space $E$, and
\[
  S_n=\PP(\Gamma^nE).
\]
For $\lambda\in E^\vee$, define divided-power contraction by
\[
 \langle\iota_\lambda f,g\rangle=\langle f,\lambda g\rangle.
\]
Then
\begin{equation}\label{eq:lift}
 g\in W_{\iota_\lambda f}\quad\Longleftrightarrow\quad
 \lambda g\in W_f.
\end{equation}
The lower witness lifts squarefreely only when $\lambda\nmid g$.
Thus the factor selected for contraction must remain as a marked
forbidden root.

\begin{definition}[Contraction and marker spaces]\label{def:polar-spaces}
Let $C_f:E^\vee\to\Gamma^{n-1}E$ be
$\lambda\mapsto\iota_\lambda f$, and put
$U_f=\PP(E^\vee)\setminus\PP(\ker C_f)$.  The first-polar rational map
is
\[
 \Phi_f:U_f\longrightarrow S_{n-1},\qquad
 [\lambda]\longmapsto[\iota_\lambda f]
.
\]
Its image is a point or a line.  For $j\geq1$, let
\[
\operatorname{Conf}_j(\PP(E^\vee))
=\{(\lambda_1,\ldots,\lambda_j):\lambda_i\ne\lambda_\ell
  \text{ for }i\ne\ell\}.
\]
The ordered marker space $U_{f,j}$ is the open subset on which every
successive contraction is nonzero.  The universal $j$-step polar flag
is the map
\[
(\lambda_1,\ldots,\lambda_j)\longmapsto
\bigl([\iota_{\lambda_1}f],\ldots,
[\iota_{\lambda_j}\cdots\iota_{\lambda_1}f]\bigr).
\]
\end{definition}

In an affine coordinate $r$ the first map is
\[
\Phi_f(r)=(a_1-ra_0,\ldots,a_n-ra_{n-1}),\qquad
\Phi_f(\infty)=(a_0,\ldots,a_{n-1}).
\]

\begin{definition}[Lower splitting package]\label{def:lower-package}
At level $(n,j)$ a lower splitting package consists of:
\begin{enumerate}[label=\textup{(\roman*)}]
\item a finite stratification
  $X_{n-j}^{\mathrm{split}}=\coprod_\alpha X_\alpha$ of the ordered
  squarefree splitting incidence over $S_{n-j}$;
\item for each $\alpha$, a specified identity or Frobenius-twisted
  geometrically integral root cover
  $Y_\alpha\to X_\alpha$, with genus bound $g_\alpha$ and point bound
  \[
       \#Y_\alpha(\F_q)\geq
       q+\kappa_\alpha-2g_\alpha\sqrt q;
  \]
\item a closed bad scheme $B_{n-j}\subset S_{n-j}$ containing the
  loci on which the chosen cover is not available;
\item branch, diagonal, old-marker, new-marker, and self-collision
  divisors whose total degree on $Y_\alpha$ is at most
  $\delta_\alpha$;
\item bounds $b$ and $c$ for the pullbacks of $B_{n-j}$ and the finite
  self-collision divisor to a noncontained first-polar line.
\end{enumerate}
The index $\alpha$ names a geometric stratum together with its constant
field and Frobenius twist.  All degrees are geometric degrees after
the displayed pullback.  ``Computed'' means that equations for
$B_{n-j}$ and every divisor are obtained by finite contraction,
elimination, factorization, and rank tests over the coefficient ring.
\end{definition}

The construction is summarized in Figure~\ref{fig:polar-flow}.  The
upper branch is the level-specific component calculation; the lower
branch is the general transverse argument.
\begin{figure}[H]
\centering
\[
\begin{array}{ccccc}
[f]\in S_n&\dashrightarrow&\Phi_f(U_f)\subset S_{n-1}
 &\longrightarrow&\text{lower ordered splitting cover}\\
&&\begin{array}{c}
\searrow\ \text{contained}\\[-2pt]
\nearrow\ \text{transverse}
\end{array}
&&\begin{array}{l}
\text{persistent/Lucas component},\\
\text{delete branch, diagonal, and markers}
\end{array}\\
&&&&\downarrow\ \text{lift by \eqref{eq:lift}}\\
&&&&\text{split squarefree witness in }W_f .
\end{array}
\]
\caption{Coherent polar induction.  The contained branch is a
level-specific component theorem; only the transverse branch is
controlled by the general point-count argument.}
\label{fig:polar-flow}
\end{figure}

\begin{theorem}[Polar-flag construction]\label{thm:polar-construction}
The maps of Definition~\ref{def:polar-spaces} commute with field
extension and the natural $\operatorname{GL}(E)$ action.  They include
the infinity chart displayed above.  If
$g\in W_{\iota_{\lambda_j}\cdots\iota_{\lambda_1}f}$, then
$\lambda_1\cdots\lambda_jg\in W_f$; the lift is squarefree exactly
when $g$ is squarefree and avoids every marker $\lambda_i$.
\end{theorem}

\begin{proof}
Contraction is defined by the perfect divided-power/symmetric-power
pairing, so it commutes with base change and change of basis.  Repeated
application of \eqref{eq:lift} gives the membership assertion.  Since
the $\lambda_i$ are distinct on the configuration space, their product
is squarefree; the product with $g$ is squarefree precisely when $g$
is squarefree and none of the $\lambda_i$ divides it.  The coordinate
formula at infinity follows by evaluating $C_f$ on the second basis
covector.
\end{proof}

\begin{definition}[Contained-component assertion]\label{def:contained}
For a specified lower package, $\mathrm{CC}(n,j)$ is the
scheme-theoretic assertion that every $f\in S_n$ for which
$\Phi_f^*B_{n-j}$ or a marker-collision section vanishes identically
belongs to the explicitly listed closed subscheme
$\cP_n\cup\cM_n$.  An identically zero collision differential is part
of $\cM_n$; it is never counted as a finite divisor.
\end{definition}

This is a level-specific theorem obligation, not an assumption hidden
inside the point-count argument.  Sections~\ref{sec:r67} and
\ref{sec:fixed} state the contained-component calculations used at
redundancies six through nine.  No assertion $\mathrm{CC}(n,j)$ is
made for an arbitrary degree.

\begin{theorem}[Effective transverse induction]\label{thm:induction}
Fix a lower splitting package as in Definition~\ref{def:lower-package}
and suppose $\mathrm{CC}(n,j)$ has been proved for that package.  Put
\[
\mathcal H_\kappa(g,\delta)=
\left\lfloor
\left(g+\sqrt{g^2+\delta-\kappa}\right)^2
\right\rfloor+1
\]
and assume $\delta\geq\kappa-g^2$.  If
\[
q\geq \max_\alpha
\mathcal H_{\kappa_\alpha}(g_\alpha,\delta_\alpha)
\quad\text{and}\quad q+1>b+c,
\]
then every split-free Hankel system belongs to the contained scheme
$\cP_n\cup\cM_n$.
\end{theorem}

\begin{proof}
For a system outside $\cP_n\cup\cM_n$, the assertion
$\mathrm{CC}(n,j)$ makes the pulled-back bad and self-collision schemes
finite of total degree at most $b+c$.  Since
$\#\PP^1(\F_q)=q+1$, choose a rational contraction parameter outside
them.  On the corresponding lower stratum the recorded point bound
gives
\[
 q+\kappa_\alpha-2g_\alpha\sqrt q>\delta_\alpha,
\]
so there is a rational point outside all deletions.  It yields a split
squarefree lower member avoiding every marker.  Equation
\eqref{eq:lift} lifts it to a split squarefree member upstairs,
contradicting split-freeness.
\end{proof}

\begin{remark}[Point-count convention]\label{rem:thresholds}
The correction $\kappa$ records the exact lower bound available for the
chosen model.  The singular cubic fiber square of
Lemma~\ref{lem:s3} uses $\kappa=0$.  The normalized genus-one covers in
the fixed-level applications use $\kappa=1$.  Thus the redundancy-eight
triple $(g,\delta,\kappa)=(1,30,1)$ gives the real cutoff
$(1+\sqrt{30})^2<42$, whose first prime power is $43$; the
redundancy-nine triple $(1,36,1)$ gives the strict cutoff $q>49$, whose
first prime power is $53$.
\end{remark}

\begin{table}[ht]
\centering
\caption{Normalized point-count thresholds used in this paper.}
\label{tab:thresholds}
\begin{tabular}{lccccc}
\toprule
Application & $g$ & $\delta$ & $\kappa$ & integer lower bound & first prime power\\
\midrule
R5 cubic fiber square & $1$ & $12$ & $0$ & $q\geq22$ & $23$\\
R6 marked cover & $1$ & $18$ & $1$ & $q\geq28$ & $29$\\
R7 marked cover & $1$ & $24$ & $1$ & $q\geq35$ & $37$\\
R8 three-marker cover & $1$ & $30$ & $1$ & $q\geq42$ & $43$\\
R9 four-marker cover & $1$ & $36$ & $1$ & $q\geq50$ & $53$\\
\bottomrule
\end{tabular}
\end{table}

\subsection{Contained-component calculations}

The ordinary rank-two carrier admits a uniform contained-line
calculation.  This does not classify other lower bad schemes, but it
closes the persistent part of every level-specific component gate.

\begin{proposition}[Contained rank-two polar lines]
\label{prop:contained-rank-two}
Let $n\geq5$, let $f\in\PP(\Gamma^nE)$, and let
\[
 A=C_f^{(2)}:\Sym^{n-2}E^\vee\longrightarrow\Gamma^2E
\]
be the three-row catalecticant.  Assume the contraction map
$E^\vee\to\Gamma^{n-1}E$ is injective, so the first-polar map is
defined on all of $\PP(E^\vee)$.  Then
\[
\Phi_f(\PP(E^\vee))\subset
\{b:\operatorname{rank}C_b^{(2)}\leq2\}
\quad\Longleftrightarrow\quad
\operatorname{rank}A\leq2
\]
scheme-theoretically and in every characteristic.  If
$\operatorname{rank}A=3$, the rank-drop parameters form a subscheme of
$\PP(E^\vee)$ of degree at most three.
\end{proposition}

\begin{proof}
For a nonzero $\lambda\in E^\vee$, multiplication identifies
$\Sym^{n-3}E^\vee$ with the hyperplane
\[
 H_\lambda=\lambda\Sym^{n-3}E^\vee\subset\Sym^{n-2}E^\vee,
\]
and functoriality gives
\[
 C_{\iota_\lambda f}^{(2)}=A|_{H_\lambda}.
\]
The forward implication is immediate.  Conversely, suppose
$\operatorname{rank}A=3$ and put $K=\ker A$, so
$\dim K=n-4$.  If the lower rank is at most two, rank--nullity on the
$(n-2)$-dimensional space $H_\lambda$ gives
\[
 \dim(K\cap H_\lambda)\geq(n-2)-2=n-4=\dim K.
\]
Thus $K\subset H_\lambda$.  Containment for every geometric
$[\lambda]$ would imply
\[
 K\subset\bigcap_{[\lambda]\in\PP(E^\vee_{\bar k})}
 \lambda\Sym^{n-3}E^\vee_{\bar k}=0,
\]
because a nonzero binary form cannot be divisible by every linear
form.  This contradicts $\dim K=n-4>0$.

The maximal minors of $C_{\iota_\lambda f}^{(2)}$ are homogeneous
cubics in $\lambda$.  Vanishing at every geometric point is therefore
identical vanishing, proving the scheme-theoretic assertion.  In the
noncontained case at least one cubic is nonzero, and the common
rank-drop scheme lies in its degree-three zero divisor.
\end{proof}

After rank one and rational split secants are removed as shallow, the
upper rank-two locus consists of the persistent tangent and
conjugate-secant carriers.  On a persistent fiber the nonsplit torus
acts by $z\mapsto z^n$, giving $T/T^n$ modulo inversion and Frobenius,
while the tangent unipotent cocycle is $z\mapsto z+nu$.

Modular carriers are equally explicit.  If $\cN$ is a nucleus in the
lower divided-power module, the complete space whose first-polar line
lies in $\cN$ is
\[
\cM_n(\cN)=
\PP\ker\left(
\Gamma^nE\longrightarrow
\operatorname{Hom}(E^\vee,\Gamma^{n-1}E/\widetilde{\cN})
\right).
\]
Lucas' theorem computes the nucleus coordinates, and this displayed
kernel computes every coherent lift without an ambient syndrome census.

\begin{remark}[Scope of the contained calculations]
The modular contraction kernel is an exact finite linear-algebra
construction.  Proposition~\ref{prop:contained-rank-two} supplies the
ordinary persistent component uniformly, but other lower bad schemes
and marker collisions still require level-specific calculations.
Neither the proposition nor Theorem~\ref{thm:induction} classifies
those additional components.
\end{remark}

\section{Redundancies six and seven: complete and gated results}
\label{sec:r67}

\subsection{Redundancy six}

For
\[
 f=(a_0:\cdots:a_5)\in\PP(\Gamma^5E),\qquad
 W_f=\ker
 \begin{pmatrix}
 a_0&a_1&a_2&a_3&a_4\\
 a_1&a_2&a_3&a_4&a_5
 \end{pmatrix},
\]
the system $W_f$ is a net of quartics off the quintic normal rational
curve.  It is split-free exactly when it has no completely split
squarefree quartic.

\begin{proposition}[Persistent net and orbit law]\label{prop:r6-persistent}
If $\gcd(W_f)=Q$ has degree two, then
\[
       W_f=Q\Sym^2E^\vee.
\]
The deep persistent locus has size $q(q+1)^2/2$.  On the tangent fiber
the unipotent coordinate is $x\mapsto x+5u$; on the irreducible
quadratic fiber the norm-one torus acts through $z\mapsto z^5$, with
inversion and coefficient Frobenius acting on $T/T^5$.
\end{proposition}

\begin{proof}
For fixed $Q$, the syndromes annihilating
$Q\Sym^3E^\vee$ form a projective line.  If $Q$ is irreducible, all
$q+1$ points remain; if $Q=\lambda^2$, one endpoint has dependent
Hankel rows, leaving $q$ points.  Hence
\[
 q(q+1)+\frac{q(q-1)}2(q+1)=\frac{q(q+1)^2}{2}.
\]
Normalize a square to $U^2$.  The nondegenerate fiber is
$[e_4+xe_5]$, and direct degree-five substitution gives
$x\mapsto x+5u$.  For an irreducible $Q$, its stabilizer identifies
the $q+1$-point fiber with
$T=\{z\in\F_{q^2}^\times:z^{q+1}=1\}$ and acts by fifth powers; the
normalizer inverts $z$, while coefficient Frobenius multiplies
$T/T^5$ by the characteristic.  This proves the stated orbit law and
the characteristic-five tangent splitting.
\end{proof}

\begin{proposition}[First-polar line and secant degree]
\label{prop:r6-secant}
For $r\in\F_q$ put
\[
 b(r)=(a_1-ra_0,\ldots,a_5-ra_4),\qquad
 b(\infty)=(a_0,\ldots,a_4).
\]
Then
\[
 (T-rU)g\in W_f\quad\Longleftrightarrow\quad g\in W_{b(r)}.
\]
If $W_f$ has trivial gcd, its polar line meets the lower secant cubic
in degree at most three.
\end{proposition}

\begin{proof}
The first equivalence follows by expanding the two Hankel equations of
$(T-rU)g$.  Let
\[
C_f=\begin{pmatrix}
a_0&a_1&a_2&a_3\\
a_1&a_2&a_3&a_4\\
a_2&a_3&a_4&a_5
\end{pmatrix}.
\]
The lower catalecticant equals
\[
C_f\begin{pmatrix}
-r&0&0\\1&-r&0\\0&1&-r\\0&0&1
\end{pmatrix}.
\]
Cauchy--Binet gives the secant determinant
\[
D_f(r)=-\Delta_{012}r^3+\Delta_{013}r^2
       -\Delta_{023}r+\Delta_{123}.
\]
It vanishes identically exactly when $\operatorname{rank}C_f\leq2$,
equivalently when $W_f$ has quadratic gcd.  In the trivial-gcd case it
therefore has at most three zeros.
\end{proof}

\begin{proposition}[Cyclic and collision degrees]\label{prop:r6-degrees}
For a trivial-gcd net:
\begin{enumerate}[label=\textup{(\roman*)}]
\item in characteristic different from two, the lower tame cyclic
carrier is a binomially rescaled Veronese surface of degree four and
contains no polar line;
\item the pointed ramification condition has degree at most six and
does not vanish identically;
\item an $S_3$ lower pencil with the marked root deleted has total
deletion degree at most eighteen.
\end{enumerate}
\end{proposition}

\begin{proof}
In syndrome coordinates the tame cyclic surface is
\[
\left\{\operatorname{diag}(1,4,6,4,1)^{-1}[Q^2]:
[Q]\in\PP(\Sym^2E^\vee)\right\}.
\]
Outside characteristic two it is the Veronese surface, of degree four,
and contains no line.  For a basis $F_0,F_1,F_2$ of the $g^2_4=W_f$,
pointed collision is the rank-one condition
\[
\operatorname{rank}
\begin{pmatrix}
F_0(r)&F_1(r)&F_2(r)\\
F_0'(r)&F_1'(r)&F_2'(r)
\end{pmatrix}\leq1.
\]
Its Wronskian divisor has degree six.  In characteristic at least five
it cannot vanish identically because the degrees are below the
characteristic.  In characteristics two and three, either all
derivatives vanish or the derivative row is a rational multiple of the
value row; the Hankel overlap excludes the first alternative, and a
gcd argument forces a common factor in the second.

Finally the R5 off-diagonal curve has the original deletion degree
twelve.  The unique cubic through a specified marker contributes at
most six ordered pairs, giving $\delta=18$ and
\[
       q-2\sqrt q>18.
\]
Its first prime-power solution is $29$.
\end{proof}

\begin{proposition}[Modular contained components]\label{prop:r6-modular}
In characteristic three the lower wild carrier is the degree-four cone
\[
       \operatorname{Join}(e_2,C_4),
\]
and no polar line of a trivial-gcd net is contained in it.  In
characteristic two the cyclic carrier is the plane
\[
       \Pi_{\mathrm{cyc}}=\PP\langle e_1,e_2,e_3\rangle,
\]
whose unique contained polar-line component is
\[
       \cN_3=\PP\langle e_2,e_3\rangle.
\]
\end{proposition}

\begin{proof}
Every line on the wild cone is a ruling.  If two consecutive Hankel
rows lie in a ruling $\langle e_2,\nu(t)\rangle$, the overlap equations
give $f=\mu\nu(t)$, including $f=\mu e_5$ at infinity.  Thus no
trivial-gcd polar line is contained.

In characteristic two, direct substitution gives
\[
g\cdot e_2=
[0:\delta\gamma:\alpha\delta+\beta\gamma:\alpha\beta:0],
\]
whose closure is $\Pi_{\mathrm{cyc}}$.  The consecutive-row overlap
places a polar line in this plane exactly when
$a_0=a_1=a_4=a_5=0$, namely when
$f\in\PP\langle e_2,e_3\rangle$.
\end{proof}

\begin{proposition}[Binary nucleus arithmetic]\label{prop:r6-nucleus}
The scheme $\cN_3(\F_q)$ is one projective-semilinear orbit of $q+1$
points.  It is split-free over $q=2^m$ exactly when $m$ is odd.
\end{proposition}

\begin{proof}
For the representative $e_3$,
\[
       W_{e_3}=\langle1,t,t^4\rangle.
\]
The root differences of $A+Bt+Ct^4$ lie in the kernel of
$u\mapsto u^4+(B/C)u$.  Four distinct rational roots exist exactly
when all roots of $u^3=B/C$ are rational, equivalently when
$3\mid q-1$, which for $q=2^m$ means $m$ is even.  The stabilizer of
$e_3$ is the Borel of order $q(q-1)$, so the orbit has $q+1$ points;
its lower-unipotent transforms and endpoint are precisely
$\cN_3(\F_q)$.
\end{proof}

The covering-radius gate is complete: Seroussi--Roth gives
$\rho(\PRS(q-5))=5$ for $q\geq9$, while Certificate R6 checks it
directly at $q=7,8$.

\begin{theorem}[Redundancy six]\label{thm:r6}
For every prime power $q\geq7$, the deep syndromes of $\PRS(q-5)$ are
the persistent stratum of size $q(q+1)^2/2$, together with:
\begin{enumerate}[label=\textup{(\roman*)}]
\item respectively $18,11,4,2,1$ exceptional $\PGL_2(q)$ orbits at
      $q=7,8,9,11,13$;
\item one characteristic-two nucleus orbit of size $q+1$ when
      $q=2^m$ and odd $m\geq5$.
\end{enumerate}
There are no other exceptional orbits.  The complete persistent orbit
law is $T/T^5$ modulo inversion and Frobenius, with a tangent
fixed/nonzero split exactly in characteristic five.
\end{theorem}

\begin{proof}
Propositions~\ref{prop:r6-persistent}--\ref{prop:r6-nucleus} classify
the common-factor and contained components and exclude a trivial-gcd
exception in the conceptual ranges: characteristic at least five from
$q\geq29$, characteristic three from its next unchecked field
$q=81$, and characteristic two outside $\cN_3$ from $q\geq32$.
Certificate R6 exhausts
\[
q\in\{7,8,9,11,13,16,17,19,23,25,27\},
\]
finds exceptional projective-orbit counts $18,11,4,2,1$ only at the
five displayed fields, and no further orbit at the remaining fields.
Together with the separate radius gate, this covers every prime power
$q\geq7$.
\end{proof}

The five small tables are exhaustive computations.  Their semilinear
orbit counts are $18,5,2,2,1$; a collision-energy invariant plus, in
odd characteristic, the quintic root-divisor type separates the
classes up to precisely Frobenius fusion.
The C498 census certificate supplies a canonical representative,
stabilizer, member histogram, and Frobenius link for every projective
orbit, while the normal-form certificate supplies the collision-energy
and root-divisor fingerprint.  Orbit counts alone are not used as
identifiers.

\subsection{Redundancy seven}

For
\[
f=(a_0:\ldots:a_6)\in\PP(\Gamma^6E),\qquad
W_f=\ker
\begin{pmatrix}
a_0&a_1&a_2&a_3&a_4&a_5\\
a_1&a_2&a_3&a_4&a_5&a_6
\end{pmatrix},
\]
the system $W_f$ is a four-dimensional web of quintics.  For an affine
marker $r$ the contractions are
\[
 b(r)=(a_1-ra_0,\ldots,a_6-ra_5).
\]
Coefficient collection gives
\[
(t-r)g\in W_f\quad\Longleftrightarrow\quad g\in W_{b(r)};
\]
the lift is squarefree exactly when $g(r)\ne0$.

\begin{proposition}[Two-marker lower cover]\label{prop:r7-pointed}
Let $x$ be a forbidden root on a trivial-gcd quartic net.  On the
generic $S_3$ bottom fiber, the diagonal and branch divisors and the
two marked-root fibers have total degree at most
\[
       12+6+6=24.
\]
Hence $q-2\sqrt q>24$, and in particular every prime power $q\geq37$,
produces a split lower quartic avoiding $x$.
\end{proposition}

\begin{proof}
The R5 off-diagonal curve is geometrically integral of arithmetic genus
one and has base deletion degree twelve.  Each marker belongs to a
unique cubic fiber, which contains at most six ordered pairs.  Delete
those two fibers and apply the $\kappa=0$ bound.
\end{proof}

\begin{proposition}[Exact-gcd-one avoidance]\label{prop:r7-gcd1}
Suppose the lower net is $\ell_xU$, where
$U=\ker\lambda\subset\Sym^3E^\vee$ has dimension three and $x\ne r$.
For $q\geq16$ it contains a split cubic avoiding both $x$ and $r$.
\end{proposition}

\begin{proof}
Put $S=\PP^1(\F_q)\setminus\{x,r\}$.  For ordered distinct
$u,v\in S$, the map
$w\mapsto\lambda(\ell_u\ell_v\ell_w)$ is linear.  If it vanishes
identically, choose any remaining $w$; otherwise it has one rational
zero.  The equations $w=x,r,u,v$ have respective bidegrees
$(1,1),(1,1),(2,1),(1,2)$ in $(u,v)$.  None is identically zero:
the first two would add a common factor, and the last two are excluded
because double-root cubics span $\Sym^3E^\vee$.  A nonzero
bidegree-$(a,b)$ divisor on $\PP^1\times\PP^1$ has at most
$(a+b)(q+1)$ rational points.  Thus the bad pairs number at most
$10(q+1)$, whereas $(q-1)(q-2)>10(q+1)$ for $q\geq16$.
\end{proof}

\begin{proposition}[Marked self-collision]\label{prop:r7-collision}
After fixed factors are removed, the self-collision divisor of the
moving $g^3_5$ has degree at most eight.
\end{proposition}

\begin{proof}
It is the ramification divisor of a separable $g^3_5$, whose degree is
$(3+1)(5-3)=8$.  If the series were inseparable, it would factor
through Frobenius and lie in the $p$-th-power subspace of quintics;
that subspace has dimension at most three, smaller than the
four-dimensional Hankel series.
\end{proof}

The characteristic-two nucleus line at redundancy six has exactly one
coherent lift, the central point $e_3$, with
\[
 W_{e_3}=\langle1,t,t^4,t^5\rangle.
\]
It is deep for odd $m$; for even $m$, the homogenization of $t^4+t$
has the five roots $\PP^1(\F_4)$.

\begin{proposition}[Central binary lift]\label{prop:r7-central}
The complete inverse image of the R6 nucleus line under consecutive
sextic contraction is $[e_3]$.  It is split-free over $\F_{2^m}$
exactly when $m$ is odd.
\end{proposition}

\begin{proof}
Putting both consecutive rows in
$\PP\langle e_2,e_3\rangle$ gives
$a_0=a_1=a_2=a_4=a_5=a_6=0$.  Thus only $e_3$ remains, with
$W_{e_3}=\langle1,t,t^4,t^5\rangle$.  For odd $m$, contraction at a
root of a hypothetical split quintic would give a split quartic on the
R6 nucleus line, contradicting Proposition~\ref{prop:r6-nucleus}.  For
even $m$, $\F_4\subset\F_q$ and the homogenization of $t^4+t$ has its
four affine $\F_4$ roots and the simple root at infinity.
\end{proof}

\begin{corollary}[Degree-six contained component]
\label{cor:r7-contained}
The assertion $\mathrm{CC}(6,1)$ for the lower rank-two carrier holds
in every characteristic.  Its contained sextics have upper
catalecticant rank at most two; after the shallow rank-one and rational
secant loci are removed, they are exactly the persistent tangent and
conjugate-secant families.
\end{corollary}

\begin{proof}
Apply Proposition~\ref{prop:contained-rank-two} with $n=6$.  The three
rational types of the quadratic recurrence are split squarefree,
repeated rational, and irreducible.  They give respectively the shallow
rational secant, tangent, and conjugate-secant loci.  The remaining
lower binary nucleus line has the separate complete lift computed in
Proposition~\ref{prop:r7-central}.
\end{proof}

\begin{theorem}[Redundancy-seven classification]
\label{thm:r7}
Certificate R7 unconditionally classifies every prime power
$q\leq32$ in its displayed domain.  For every prime power $q\geq13$,
the split-free syndromes
at redundancy seven are the persistent stratum, together with the
central nucleus point when $q=2^m$ and $m$ is odd.  At $q=7,8,9,11$ the additional
split-free orbit-size profiles are:
\[
\begin{array}{c|l}
7&56,\ 84^5,\ 112^2,\ 168^{45},\ 336^{141}\\
8&63,\ 72,\ 84^3,\ 168^4,\ 252^{24},\ 504^{86}\\
9&180^3,\ 240^6,\ 360^{18},\ 720^{27}\\
11&264^2,\ 440,\ 660^2 .
\end{array}
\]
This is the complete all-field split-free classification.  Since
Seroussi--Roth gives $\rho(\PRS(q-6))=6$ for every prime power
$q\geq11$, it is also the complete deep-hole classification in that
range.  At $q=7,8,9$ the table is not promoted to a code deep-hole
classification without a separate covering-radius calculation.  The
persistent orbit law is $T/T^6$ modulo inversion and Frobenius, with
tangent splitting in characteristics two and three.
\end{theorem}

\begin{proof}
For $q\geq37$, Propositions~\ref{prop:r7-pointed}--\ref{prop:r7-collision}
give the transverse witness after at most three catalecticant
intersections, eight self-collisions, and, in characteristic two, one
intersection with the nucleus line.
Corollary~\ref{cor:r7-contained} leaves exactly the persistent locus and
Proposition~\ref{prop:r7-central}.  Certificate R7 covers
\[
q\in\{7,8,9,11,13,16,17,19,23,25,27,29,31,32\}
\]
by an independent five-secant replay.  The radius theorem promotes the
split-free list only from $q\geq11$.
\end{proof}

For the four finite rows the C509 certificate contains the canonical
representative, stabilizer, persistent/central flags, and Frobenius link
of every orbit.  The displayed size profile is only a compact summary;
repeated sizes are distinguished by those canonical records.

The $q=19$ lower net $\langle1,t^3,t^4\rangle$ has six split members,
all containing the marked point at infinity, yet no coherent sextic
polar line lies in its orbit.  This is the concrete reason
Theorem~\ref{thm:induction} is formulated for parameterized flags rather
than a set of bad lower fibers.

\begin{table}[ht]
\centering
\small
\caption{Field-range closure ledger.  ``Certificate'' means a complete
finite classification with canonical representatives, not merely an
orbit count.}
\label{tab:field-ranges}
\begin{tabular}{lll}
\toprule
Result & prime powers & closure mechanism\\
\midrule
R5 & $7,8,9,11,13,16,17,19$ & direct census and independent replay\\
   & $q\geq23$ & cubic-cover geometry, Lemma~\ref{lem:s3}\\
R6 & $7,8,9,11,13,16$ & direct syndrome/radius census\\
   & $17,19,23,25,27$ & finite structural bridge certificate\\
   & $q\geq29$ & marked-cover point bound and contained components\\
R7 & $7,8,9,11$ & direct split-free census\\
   & $13,16,17,19,23,25,27,29,31,32$ &
       coherent-polar finite bridge certificate\\
   & $q\geq37$ & marked-cover point bound and contained components\\
\bottomrule
\end{tabular}
\end{table}

The R7 rows classify split-free syndromes.  Conversion to code
deep-hole directions uses the separate covering-radius gate and is
made only for $q\geq11$.  The public certificate schema and exact
replay boundaries are given in the supplement; the table does not
promote orbit-size profiles to identifiers.

\section{Fixed-level high-field theorems}\label{sec:fixed}

\subsection{Redundancy eight}

\begin{proposition}[Three-marker bottom cover]\label{prop:r8-bottom}
After three contractions, the generic bottom object is the R5
off-diagonal $(2,2)$ curve.  The diagonal and branch deletion costs
twelve and three marked roots cost eighteen, so
$\delta=30$.  The inequality
\[
       q+1-2\sqrt q>30
\]
holds at every prime power $q\geq43$.  The first-polar line has
transverse/collision budget
\[
       3+1+10=14,
\]
where ten is the ramification degree of a moving $g^4_6$.
\end{proposition}

\begin{proof}
Markers do not change the bottom curve or its geometric monodromy;
each removes at most the six ordered pairs in its unique cubic fiber.
The point bound is the $\kappa=1$ row of
Table~\ref{tab:thresholds}.  A noncontained line meets the lower
rank-two catalecticant carrier in degree at most three, and the lower
binary central point in at most one parameter.  After fixed factors are
removed, self-collision is ramification of a separable $g^4_6$, of
degree $(4+1)(6-4)=10$.  Separability follows because the
$p$-th-power subspace has dimension smaller than the five-dimensional
series.
\end{proof}

\begin{proposition}[Modular lifts at degree seven]\label{prop:r8-modular}
The consecutive-row Lucas calculation leaves no nonzero lift in
characteristic two, the line $\PP\langle e_2,e_5\rangle$ in
characteristic three, and the line
$\PP\langle e_3,e_4\rangle$ in characteristic five.  The latter two
lines are shallow in the range $q\geq43$.
\end{proposition}

\begin{proof}
Lucas' criterion applied to the consecutive binomial rows gives the
three stated supports; the exact row reductions are replayed by
Certificate R8.  In characteristic five, the homogenization of
$t^5-t$ supplies six distinct projective roots and satisfies the two
required middle-coefficient equations throughout the lift.

In characteristic three, $\PGL_2(q)$ is transitive on the rational
points of $\PP\langle e_2,e_5\rangle$, so it suffices to treat $e_2$.
Choose nonzero $a,b,c$ with no two equal up to sign and
$a^2+b^2+c^2=0$.  The nonsingular conic has $q+1$ points; deleting the
three coordinate and six equality/antiequality sections costs at most
eighteen, so a choice exists for $q\geq27$.  The reciprocal pairs
$\pm a^{-1},\pm b^{-1},\pm c^{-1}$ give six distinct roots and the two
required elementary-symmetric coefficients vanish.
\end{proof}

\begin{remark}[R8 contained gate]\label{rem:r8-contained}
Proposition~\ref{prop:contained-rank-two} proves
scheme-theoretically that a degree-seven first-polar line contained in
the lower rank-two carrier is itself on the upper rank-two carrier.
Thus the ordinary persistent identity is no longer a gate.  What
remains unproved is that the other lower bad strata and the third
marker create no additional contained component.
Proposition~\ref{prop:r8-bottom} proves the transverse calculation, not
that residual assertion; we retain it as the reduced
$\mathrm{CC}(7,1)$ gate.
\end{remark}

\begin{theorem}[Conditional redundancy eight]\label{thm:r8}
Assume $\mathrm{CC}(7,1)$.  For every prime power $q\geq43$, the deep
syndromes of $\PRS(q-7)$ are
exactly the persistent tangent and conjugate-secant families, of total
size $q(q+1)^2/2$.  Their orbit law is $T/T^7$ modulo inversion and
Frobenius.  The tangent family splits exactly in characteristic seven.
\end{theorem}

\begin{proof}
Outside the contained locus, Proposition~\ref{prop:r8-bottom} chooses a
good first contraction and supplies a three-marker lower witness.
Proposition~\ref{prop:r8-modular} eliminates every modular lift.  The
assumed contained statement leaves only the persistent families, and
the imported covering-radius theorem promotes split-free directions to
deep holes.
\end{proof}

This theorem makes no assertion about additional deep orbits for
$q<43$.

\subsection{Redundancy nine}

\begin{proposition}[Four-marker budgets]\label{prop:r9-budgets}
Four markers give deletion degree
\[
       12+4\cdot6=36,
\]
and the first prime-power solution of
$q+1-2\sqrt q>36$ is $53$.  The transverse/collision budget is at most
\[
       3+2+12=17.
\]
\end{proposition}

\begin{proof}
The bottom curve contributes twelve diagonal/branch points and each
marker its six ordered pairs.  The first-polar line meets the
persistent catalecticant carrier in degree at most three and the
linear nucleus support in at most two parameters.  Self-collision is
ramification of a separable $g^5_7$, of degree
$(5+1)(7-5)=12$.  The normalized point bound is the
$(g,\delta,\kappa)=(1,36,1)$ row of
Table~\ref{tab:thresholds}.
\end{proof}

The only substantial modular carrier occurs in characteristic seven
and is a projective binary-quartic space.  Fixing a split quintic
$P(t)=\sum_{i=0}^5p_it^i$ and a carrier point
$h=(a_0,\ldots,a_4)$, put
\[
 H_j=\sum_{i=0}^4a_ip_{i+j}
\]
with $p_k=0$ outside $0\leq k\leq5$, and
\[
 D=H_0H_2-H_1^2,\quad
 N_s=H_{-1}H_2-H_0H_1,\quad
 N_u=H_{-1}H_1-H_0^2.
\]
Off $D=0$, the unique residual quadratic is
\[
       Dt^2-N_st+N_u,
\]
with branch polynomial $K=N_s^2-4N_uD$.  The exact residual identities
are proved next; the geometric exhaustion of transported slices is
kept as a separate printed gate.

\begin{proposition}[Residual quadratic identities]\label{prop:r9-residual}
With $p_k=0$ outside $0\leq k\leq5$, the two consecutive Hankel
equations for
\[
       P(t)(t^2-st+u)
\]
are
\[
H_0-sH_1+uH_2=0,\qquad
H_{-1}-sH_0+uH_1=0.
\]
On $D\ne0$, Cramer's rule gives
\[
       s=N_s/D,\qquad u=N_u/D.
\]
The branch divisor is $K=0$, fixed/residual collision is
$\operatorname{Res}(P,Dt^2-N_st+N_u)=0$, and fixed-root collision is
$\operatorname{Disc}(P)=0$.
\end{proposition}

\begin{proof}
Expand the two consecutive coefficients annihilated by the carrier
quartic.  The displayed linear system has determinant $D$, so Cramer's
rule gives the residual factor.  Its discriminant is $K$.  The two
resultants have their usual intrinsic collision meanings; consequently
the construction commutes with scalar, $\PGL_2$, and Frobenius
transport.
\end{proof}

\begin{proposition}[One-moving-root slice]\label{prop:r9-slice}
Let $P=R(t)(t-x)$ with four fixed distinct roots.  Then
$D,N_s,N_u$ have degree at most two in $x$, and $K$ has degree at most
four.  If the squarefree reduction of $K$ is nonconstant and not a
square, the normalization of $y^2=K(x)$ is geometrically integral of
genus at most one.
\end{proposition}

\begin{proof}
Each $H_j$ is affine in $x$, so its $2\times2$ minors are quadratic and
$K$ is quartic.  In odd characteristic the squarefree nonsquare
equation defines a geometrically integral double cover; the
hyperelliptic genus formula gives
$\lfloor(\deg K_{\mathrm{red}}-1)/2\rfloor\leq1$.
\end{proof}

\begin{remark}[Six-slice and base-selection gates]\label{rem:r9-gates}
Certificate R9 records six reduced discriminant polynomials for the
squarefree quartic normal form, their unit gcd, and four multiple-root
controls.  The development report also states a nonzero rational
base-selection polynomial of individual degree at most $24$ and total
degree at most $96$.  A refereeable proof must print the transporter,
the six polynomials and a B\'ezout identity, the four reduced branch
polynomials, and the nonzero base polynomial with its degree accounting.
Until those data appear in the paper or immutable supplement, the
six-slice component theorem and the $q\geq343$ base selection are
explicit gates rather than consequences of the certificate.
\end{remark}

\begin{proposition}[Good-slice deletion]\label{prop:r9-deletion}
Once a rational base with an integral slice is fixed, the moving/fixed
diagonal, determinant, branch, four fixed-root collisions, and
moving-root collision have degrees
\[
       4,\ 2,\ 4,\ 8,\ 4,
\]
respectively.  Their total is $22$, so
$q+1-2\sqrt q>22$ supplies a split witness.
\end{proposition}

\begin{proof}
The first three degrees follow from
Proposition~\ref{prop:r9-slice}; each fixed-root collision is linear in
the moving parameter on each of four residual sheets, and the final
resultant has degree four.  Removing their union from the normalized
genus-at-most-one slice proves the claim.
\end{proof}

\begin{proposition}[Characteristic-seven finite bridges]
\label{prop:r9-char7}
At $q=7$, the formal carrier has exactly $819$ projective rootless
quartics.  At $q=49$, Certificate R9-49 exhausts all
\[
       \frac{49^5-1}{48}=5{,}884{,}901
\]
projective carrier points and finds a split witness for each.
\end{proposition}

\begin{proof}
At $q=7$, a split squarefree degree-seven projective divisor is the
complement of one of the eight rational points.  Inclusion--exclusion
gives
\[
\frac{7^5-1}{6}-8\frac{7^4-1}{6}
+28\frac{7^3-1}{6}-56\frac{7^2-1}{6}
+70\frac{7-1}{6}=819.
\]
The $q=49$ clause is deliberately certificate-only: one lexicographic
five-root search stops after every carrier point has a witness.  The
residual algebra is independently replayed, but there is no second
exhaustive carrier implementation.
\end{proof}

\begin{proposition}[R9 persistent orbit law]\label{prop:r9-orbits}
Put $d=\gcd(8,q+1)$.  The sigma-fiber projective orbits are the
inversion orbits in $C_d=T/T^8$, and coefficient Frobenius acts by
multiplication by the characteristic.  The tangent fiber is transitive
unless the characteristic is two, when its zero and nonzero parts are
separate.
\end{proposition}

\begin{proof}
The nonsplit torus acts by eighth powers and its normalizer inverts the
torus coordinate.  Thus inversion-fixed and paired classes have
stabilizers $2d$ and $d$, respectively; coefficient Frobenius is the
$p$-power map.  For $d=8$, it fuses the two odd pairs exactly when
$p\equiv\pm3\pmod8$.  On the tangent fiber the unipotent action is
$z\mapsto z+8u$, which vanishes exactly in characteristic two.
\end{proof}

\begin{proposition}[Other modular lifts]\label{prop:r9-other-modular}
The consecutive Lucas supports leave a line in characteristic five and
the quartic carrier above in characteristic seven; no other
characteristic has a nonzero coherent lift.  The former line is shallow
for every $q>5$.
\end{proposition}

\begin{proof}
Lucas' theorem and the consecutive-row overlap give the two supports.
In characteristic five, homogenize $t^5-t$ to include its simple root
at infinity and multiply by a linear factor $t-a$ with
$a\notin\F_5$.  The seven roots are distinct and the two required
Hankel coefficients vanish.  Such an $a$ exists for every $q>5$.
\end{proof}

\begin{theorem}[Conditional redundancy nine]\label{thm:r9}
Assume the six-slice and rational-base assertions of
Remark~\ref{rem:r9-gates}, together with the degree-eight
contained-component assertion.  For every prime power $q\geq53$, the
deep syndromes of $\PRS(q-8)$ are
exactly the persistent tangent and conjugate-secant families, of total
size $q(q+1)^2/2$.  Their orbit law is $T/T^8$ modulo inversion and
Frobenius; the tangent family splits exactly in characteristic two.
\end{theorem}

\begin{proof}
Proposition~\ref{prop:r9-budgets} handles every transverse
nonmodular point.  The assumed contained assertion leaves the
persistent and modular supports.  Proposition
\ref{prop:r9-other-modular} eliminates characteristic five.  For
characteristic seven, the only powers below $343$ are $7$ and $49$,
both below the headline range; at every $q=7^m\geq343$, the assumed
base-selection assertion and Proposition~\ref{prop:r9-deletion}
produce a witness.  The covering-radius theorem completes the
promotion to deep holes.
\end{proof}

Proposition~\ref{prop:r9-char7} is a carrier boundary, not part of the
headline classification: the $q=7$ count is not a deep-hole statement
for the inadmissible parameter $\PRS(-1)$, and the $q=49$ pass is not
an ambient $\PP^8$ census.

\section{Characteristic two: the ordered Hessian}\label{sec:hessian}

In characteristic two the divided binary-cubic discriminant becomes
the doubled quadric $(AD+BC)^2$, so its pullback cannot distinguish
generic splitting.  The correct replacement retains an ordered
residual root.

Let $\ell=\PP\langle U,V\rangle\subset\PP(\Gamma^3E)$ and write
\[
\Hess^{\mathrm{div}}(uU+vV)
 =N_u(u,v)X^2+N_s(u,v)XY+D(u,v)Y^2.
\]
The ordered-Hessian incidence
\begin{equation}\label{eq:hessian}
N_u(u,v)X^2+N_s(u,v)XY+D(u,v)Y^2=0
\end{equation}
has bidegree $(2,2)$ and arithmetic genus one.  Off $N_s=0$, its
Artin--Schreier class is $DN_u/N_s^2$.

The parameter space of cubic lines is
$G(1,\Gamma^3E)$.  Let
\[
\Sigma=\{\PP(QE):[Q]\in\PP(\Sym^2E)\}\subset G(1,\Gamma^3E)
\]
be the Veronese surface of common-quadratic lines.  In the ambient
cubic space, let $\mathcal Q_0=\{AD+BC=0\}\subset\PP(\Gamma^3E)$.
Its two families of lines give two Fano ruling conics in the
Grassmannian; one is the tangent-conic boundary of $\Sigma$, and we
call the other $\cN$.  A root-compatible pullback is the family of
lines obtained by fixing a split factor $R$ and varying the final two
linear factors in the consecutive Hankel contraction of a syndrome.

\begin{figure}[H]
\centering
\[
\begin{array}{ccccc}
G(1,\Gamma^3E)
&\supset&
\text{degenerate ordered-Hessian locus}
&=&
\Sigma\ \cup\ \operatorname{Fano}_1(\mathcal Q_0)\\[4pt]
&&\downarrow\ \text{root-compatible pullback}&&\\[4pt]
&&
\begin{array}{c}
\text{persistent component}\\
\cup\\
\text{Lucas-nucleus component}
\end{array}
&&
\cN:\ \text{no rank-two pullback}.
\end{array}
\]
\caption{Ordered-Hessian geometry.  The degeneracy locus is computed
in the Grassmannian first; root compatibility then selects the
contained syndrome carriers and eliminates the complementary ruling.}
\label{fig:hessian-flow}
\end{figure}

\begin{theorem}[Ordered-Hessian geometric degeneracy]
\label{thm:hessian}
After removing the forced vertical factors from intersections with the
twisted cubic, the complete geometric degeneracy locus of
\eqref{eq:hessian} consists of:
\begin{enumerate}[label=\textup{(\roman*)}]
\item the Veronese surface $\Sigma$ of lines $\PP(QE)$, on which the incidence is
      separably reducible; and
\item the two Fano ruling conics of $\mathcal Q_0$, on which the root
      projection of \eqref{eq:hessian} is inseparable.
\end{enumerate}
The complementary ruling $\cN$ has no nontrivial rank-two root-compatible
Hankel pullback.
\end{theorem}

\begin{proof}
Once vertical and horizontal factors are removed, any separable
factorization of a $(2,2)$ curve must be
$(1,1)+(1,1)$.  The two factors are graphs of projectivities.
Semisimple and unipotent normal forms show that only the lines $QE$
occur.  The ambient equation $AD+BC=0$ gives two quadric rulings.
Substituting
the consecutive Hankel contraction forms into the complementary ruling
forces overlapping triples of linear forms to be proportional, so the
root-compatible line has rank at most one.
\end{proof}

\begin{remark}[Ordered-Hessian pullback gates]\label{rem:hessian-gates}
Two further statements are required for the advertised PRS containment
corollary.  First, the pullback of $\Sigma$ must be proved
scheme-theoretically to equal the persistent catalecticant/Lucas union;
this is a contained-component theorem, not a consequence of
Theorem~\ref{thm:hessian}.  Second, effective base selection must use a
single nonzero bad polynomial whose zero locus contains the entire
union of $\Sigma$, both rulings, and the vertical/Chow loci.  The report
bounds each individual generator by degree four in each selected root,
but does not show that one polynomial with the same bound contains the
union.  Multiplying component equations changes the stated threshold.
These are explicit proof gates.
\end{remark}

\begin{proposition}[Conditional effective Hessian corollary]
\label{prop:hessian-effective}
Assume the two gates of Remark~\ref{rem:hessian-gates}, with a global
bad polynomial of individual degree at most four.  For syndrome degree
$n\geq5$, outside the persistent/Lucas union an integral reduced slice
exists whenever
\[
q\geq
\min\left\{\frac{(n-4)(n+3)}2+1,\ 5(n-4)\right\}.
\]
Its deletion degree is at most $3n-4$.  If also
$q+1-2\sqrt q>3n-4$, every split-free syndrome lies in that union.
\end{proposition}

\begin{proof}
For a split base $R$ of degree $n-4$, the root-compatible Pl\"ucker
coordinates are quadratic in its coefficients.  The assumed global bad
polynomial has individual root degree at most four.
Schwartz--Zippel after adjoining the Vandermonde gives the first bound;
evaluation on pairwise disjoint five-element blocks gives the second.
The moving diagonal, determinant, branch, fixed-root collision, and
final collision divisors have degrees
\[
       n-4,\ 2,\ 2,\ 2(n-4),\ 4,
\]
whose sum is $3n-4$.  The normalized integral $(2,2)$ slice has genus at
most one, so the point bound leaves a rational ordered residual root
outside every deletion, and contraction lifts it squarefreely.
\end{proof}

Theorem~\ref{thm:hessian} is a containment theorem.  It does not assert
that every Lucas-carrier point is deep; that remains an arithmetic
question for each level.

\section{Lucas carriers and their arithmetic}\label{sec:lucas}

Let $d=2^s$ with $s\geq2$ in characteristic two.  Lucas' theorem gives the top
nucleus of the degree-$d$ normal rational curve as
\[
 \PP\langle e_1,\ldots,e_{d-1}\rangle,
\]
whose coherent degree-$(d+1)$ lift is
\[
 \cM_{d+1}=\PP\langle e_2,\ldots,e_{d-1}\rangle.
\]
At the distinguished Borel endpoint $e_{d-1}$ the Hankel kernel
contains
\[
       \cU_s=\langle1,t,t^d\rangle.
\]

\begin{proposition}[Lucas overlap kernel]\label{prop:lucas-kernel}
For $d=2^s$, the lower nucleus and its coherent lift are exactly
\[
\PP\langle e_1,\ldots,e_{d-1}\rangle,\qquad
\PP\langle e_2,\ldots,e_{d-1}\rangle.
\]
At $e_{d-1}$, $\cU_s$ is endpoint-specific; the intersection common to
the full carrier is only $\langle1,t^d\rangle$.
\end{proposition}

\begin{proof}
Lucas' theorem gives
$\binom di\equiv0\pmod2$ for $0<i<d$.  Applying this support condition
to the two consecutive contraction rows removes the two endpoint
coordinates and gives the displayed lift.  For a basis point $e_i$,
the Hankel equations kill the adjacent coefficients $b_{i-1},b_i$.
At $i=d-1$, the monomials $1,t,t^d$ survive.  Varying $i$ over the
carrier leaves only $1,t^d$ in the intersection.
\end{proof}

\begin{proposition}[Linearized root cover]\label{prop:linearized}
On the generic locus $BC\ne0$, the normalized ordered-root cover of
$A+Bt+Ct^d$ has geometric
monodromy $\operatorname{AGL}_1(\F_d)$, minimal constant field
$\F_d$, and based deck group $C_{d-1}$.  Coefficientwise Frobenius acts
on the based group by multiplication by two, with exact order $s$.
The net contains a split squarefree member over $\F_{2^m}$ if and only
if $s\mid m$.
\end{proposition}

\begin{proof}
After dividing by $C$, choose
\[
 \beta^{d-1}=B/C,\qquad y^d+y=A/(C\beta^d).
\]
The roots are $\{\beta(y+c):c\in\F_d\}$.  Their affine symmetries give
$\operatorname{AGL}_1(\F_d)$.  The Kummer extension has degree $d-1$
because the valuation of $B/C$ at its simple zero is one.  After
adjoining $\beta$, the right side of the Artin--Schreier equation has a
simple pole in $A/C$; an element of the form $z^d+z$ has pole order
divisible by $d$.  Hence the additive extension has degree $d$, the
cover is connected, and its geometric monodromy is the full affine
group.

Ratios of nonzero root differences recover $\F_d$, so this is the
minimal constant field.  Boundary members with $C=0$ have degree at
most one, while those with $B=0$ have zero derivative; neither is split
squarefree of degree $d$.  Thus a split member forces
$\F_d\subset\F_{2^m}$, equivalently $s\mid m$.  Conversely $t^d+t$
splits simply under that inclusion.  The equality
$\operatorname{ord}_{2^s-1}(2)=s$ follows from
\[
\gcd(2^r-1,2^s-1)=2^{\gcd(r,s)}-1.
\]
\end{proof}

At $s=2$ this recovers the odd/even extension law at redundancies six
and seven.  The first fresh carrier is $s=3$, in syndrome degree nine.
The $\F_8$-linear section has order-three component Frobenius, but its
four quotient directions change the full answer.

\begin{proposition}[The $e_7$ orbit]\label{prop:e7-orbit}
The stabilizer of $e_7\in\PP(\Gamma^9E)$ is the upper Borel, so its
$\PGL_2$-orbit is $\PP^1$.  At the distinguished point,
\[
W_{e_7}=\langle1,t,t^2,t^3,t^4,t^5,t^8\rangle.
\]
\end{proposition}

\begin{proof}
Lucas parity and extraction of the $x^2y^7$ coefficient give
\[
g e_7=(ad+bc)^2
\bigl(b^5e_2+b^4d\,e_3+bd^4e_6+d^5e_7\bigr).
\]
This is proportional to $e_7$ exactly when $b=0$.  The apolar action is
contragredient, so it transports both $W_f$ and split root divisors.
The displayed kernel follows from the two Hankel equations.
\end{proof}

\begin{proposition}[Framed Artin--Schreier quotient]
\label{prop:e7-as}
Order roots $0,1,a,b,c,d,u,v$ and put
\[
S=a+b+c+d,\quad E=ab+ac+ad+bc+bd+cd,\quad
h=1+S,\quad p=1+E+S+S^2.
\]
On the open set where the six fixed roots are distinct,
$h,p,1+h+p$ are nonzero, and
$x^2+hx+p\ne0$ for $x=a,b,c,d$, the remaining pair is the
geometrically integral Artin--Schreier cover
\[
       y^2+y=p/h^2.
\]
Its rational lifting law is
$\Tr_{\F_{2^m}/\F_2}(p/h^2)=0$.
\end{proposition}

\begin{proof}
The equations $e_1=e_2=0$ give $u+v=h$ and $uv=p$.  On $h\ne0$, put
$y=u/h$.  Along $h=0$, the leading coefficient of $p/h^2$ is
\[
D=1+bc+bd+cd+B+B^2,
\]
and $\partial D/\partial b=1+c+d\ne0$, so $D$ is not a square in the
residue field.  If $g^2+g$ has a pole of order two, its leading
coefficient is the square of the principal coefficient of $g$.
Therefore the Artin--Schreier class remains nonzero after extending
constants, proving geometric integrality.  The trace criterion gives
the rational lifting law.
\end{proof}

\begin{proposition}[Additive frame cover]\label{prop:e7-additive}
The family
\[
t^8+A_4t^4+A_2t^2+A_1t+A_0,\qquad A_1\ne0,
\]
has a connected affine-frame cover with geometric monodromy
$\operatorname{AGL}_3(\F_2)$.
\end{proposition}

\begin{proof}
For an $\F_2$-space $U$ and $w\notin U$,
\[
L_{U+\langle w\rangle}=L_U^2+L_U(w)L_U.
\]
Induction gives exponents $8,4,2,1$ and
$A_1=\prod_{v\ne0}v\ne0$.  An affine frame is an origin plus three
independent differences.  The coefficient map is generically finite
with exactly
$8|\operatorname{GL}_3(\F_2)|=1344$ frames in a fiber.  Its frame
function field therefore has degree $1344$ over the coefficient field
and invariant field equal to that coefficient field.
\end{proof}

\begin{theorem}[Shallowness of the degree-nine $e_7$ orbit]
\label{thm:e7}
For every $m\geq3$, the entire $\PGL_2$-orbit of $e_7$ is shallow over
$\F_{2^m}$.  The number of monic additive-affine witnesses is
\[
\frac{q(q-1)(q-2)(q-4)}{1344}.
\]
\end{theorem}

\begin{proof}
Every three-dimensional $\F_2$-subspace $V\subset\F_{2^m}$ gives
\[
L_V(t)=\prod_{v\in V}(t-v)
      =t^8+A_4t^4+A_2t^2+A_1t.
\]
The $t^7,t^6$ coefficients vanish, so $L_V\in W_{e_7}$, and
$A_1\ne0$ makes its roots simple.  Every admissible binary
redundancy-ten field has $m\geq4$.  Proposition~\ref{prop:e7-orbit}
transports the witness around the orbit.

There are
\[
\binom m3_2=\frac{(q-1)(q-2)(q-4)}{168}
\]
three-spaces, each with $q/8$ affine translates.  Since
$L_{a+V}(t)=L_V(t)+L_V(a)$ and the affine root set recovers its
difference space and coset, there is no overcount.
\end{proof}

Thus the order-three law of Proposition~\ref{prop:linearized} belongs
to a proper $\F_8$-linear section, not to deepness of the full
$e_7$ kernel.  The other intrinsic strata of the six-dimensional
degree-nine carrier are not classified here.

\section{Verification and formal boundary}\label{sec:verification}

The conceptual theorems above are accompanied by deterministic,
committed evidence bundles.  Table~\ref{tab:evidence} records which
claims use exhaustive computation.  Every named bundle contains the
primary generator, a compact JSON or text certificate, an independent
replay where stated, and a SHA-256 manifest.
The load-bearing generator, certificate, and replay hashes were
rechecked during manuscript integration.  The C491 and C498 manifests
also include their living theorem reports; those two report-only hashes
became stale after later closeout expansions.  No computational
artifact hash is stale, and no predecessor manifest is silently
rewritten here.  The public names Certificate R5 and Certificate R6
refer to those bundles through the release-level schema; the C-numbers
are retained only as provenance.

\begin{figure}[H]
\centering
\[
\begin{array}{ccc}
&\text{manuscript proof}\ \downarrow&\\[2pt]
\text{classical import}
&\longrightarrow\ \boxed{\text{paper claim}}\ \longleftarrow&
\text{certificate/replay}\\[2pt]
&\text{kernel implication}\ \uparrow&
\end{array}
\]
\caption{Claim-level trust routes.  A mixed theorem may use several
routes, but each route keeps its own boundary: certificates close
finite domains, imports retain their hypotheses, and Lean proves only
the printed implication.}
\label{fig:trust-routes}
\end{figure}

\small
\begin{longtable}{>{\raggedright\arraybackslash}p{0.19\textwidth}
                  >{\raggedright\arraybackslash}p{0.25\textwidth}
                  >{\raggedright\arraybackslash}p{0.47\textwidth}}
\caption{Public verification map.  Exact filenames, commands, hashes,
domains, and stop conditions belong to the electronic supplement.}
\label{tab:evidence}\\
\toprule
Public label & Verified domain & Replay and trust boundary\\
\midrule
\endfirsthead
\toprule
Public label & Verified domain & Replay and trust boundary\\
\midrule
\endhead
\bottomrule
\endfoot
Certificate R5 & nineteen finite fields & independent pointwise and
orbitwise replay; canonical representatives and exhaustion records\\
Certificate R6 & direct scan through $q=16$; finite bridge below $29$ &
independent syndrome replay and separate radius gate\\
Certificate R6-NF & small exceptional normal forms & deterministic
normal-form reconstruction; not a second implementation\\
Certificate R7 & $q=7,8,9,11$ and finite bridge below $37$ &
independent five-secant/orbit checks; split-free and radius flags remain
separate\\
Certificate R8 & fixed-level characteristic cases and bounds &
independent algebra/nucleus replay; not an ambient census\\
Certificate R9 & residual algebra and transported slices & independent
residual/slice replay; geometric integrality remains a paper proof\\
Certificate R9-49 & characteristic-seven carrier at $q=49$ & one
exhaustive carrier implementation; not a whole-code census\\
Certificate Hessian & bounded ordered-Hessian algebra & regression
only; does not replace the geometric component proof\\
Certificate Lucas & recorded Lucas parameter domain & independent
arithmetic replay\\
Certificate e7 & quotient-cover open set and additive specialization &
independent quotient-cover replay\\
\end{longtable}
\normalsize

The public labels are defined by
\texttt{supplement/CERTIFICATE-SCHEMA.md}.  Literal development commands
and working directories are recorded in
\texttt{supplement/REPRODUCING.md} and
\texttt{verification-map.md}.  The release manifest is intentionally
marked incomplete until it contains a repository URL, immutable tag and
commit, archive identifier, DOI, byte counts, and hashes.  Consequently
the present development draft is not an externally archived
proof-complete release.

The redundancy-nine residual algebra and synthesis implication are
also formalized in Lean.  The import gate and its axiom-audit module are
\[
\begin{split}
&\texttt{RelativeConicArcs.Gates.PRSRedundancyNine},\\
&\texttt{RelativeConicArcs.Gates.PRSRedundancyNineAxiomAudit}.
\end{split}
\]
Lean checks divided-power contraction, both residual Hankel identities,
the determinant and branch formulas, the residual roots and their
sum/product, the witness-to-shallow implication, the persistent-family
cardinality, and the orbit-count table from explicit arithmetic-case
and family-count inputs.  It does not derive that table from a formal
$\PGL_2/\PGammaL_2$ action.  Geometric integrality,
Hasse--Weil existence after deletions, identification of geometric
families with PRS syndromes, and exhaustion remain explicit structure
fields rather than hidden axioms.  The audit reports only the standard
\texttt{propext}, \texttt{Classical.choice}, and \texttt{Quot.sound}
dependencies.

\begin{table}[ht]
\centering
\small
\caption{Separation of mathematical input from kernel-checked
formalization.}
\label{tab:formal-trust}
\begin{tabular}{p{0.19\textwidth}p{0.34\textwidth}p{0.37\textwidth}}
\toprule
Result & Mathematical input & Kernel-checked content\\
\midrule
R9 residual synthesis &
six-slice component theorem, genus/deletion bounds, coding
identification, and exhaustion &
divided-power contraction, residual identities, determinant/branch
formulas, roots, and witness-to-shallow implication\\
R9 family inventory &
the actual $\PGL_2/\PGammaL_2$ action and geometric family
identification &
arithmetic case inputs, family-count inputs, and the resulting numerical
table\\
\bottomrule
\end{tabular}
\end{table}

\section{Provenance and responsibility}\label{sec:provenance}

This research was human-directed and agent-generated.  The author
directed the programme and takes responsibility for the manuscript; AI
systems are not authors.  Section~\ref{sec:verification} assigns every
published claim an exact trust route.  The corresponding source,
certificates, replays, and pinned Lean targets will be archived at the
paper-specific export repository and versioned DOI recorded in the
release manifest; the development monorepo is not a publication
artifact.  The
specific model and tool disclosure will follow the policy of the
selected repository and journal.  No placeholder archive field is
presented here as a completed public release.

\section{Exact boundary and further questions}\label{sec:boundary}

The results proved here do not constitute an arbitrary-redundancy
classification.
\begin{enumerate}[leftmargin=2em]
\item Redundancy five is complete for every $q\geq7$.
\item Redundancy six is a complete all-field deep-hole classification.
      Redundancy seven has a complete all-field split-free
      classification; its $q=7,8,9$ covering-radius gate remains
      separate.
\item Redundancies eight and nine have complete transverse arguments
      in the displayed high-field ranges, but their contained and
      slice/base-selection gates remain explicit.  No unconditional or
      bounded-field completion is inferred.
\item The characteristic-two ordered-Hessian theorem classifies the
      geometric degeneracy strata.  Its root-compatible pullback and
      effective base-selection constant retain the stated gates, and
      it does not classify arithmetic deepness on every Lucas carrier.
\item The $e_7$ orbit is only one intrinsic stratum of the degree-nine
      Lucas carrier.  The other strata, and therefore a redundancy-ten
      synthesis, remain open.
\item The effective polar-induction theorem requires a new
      level-specific monodromy/component calculation at each
      redundancy.  A field census cannot replace this input.
\end{enumerate}

The structural conclusion is nevertheless stable.  Beyond redundancy
four, deepness is no longer controlled by one invariant or by isolated
bad fibers.  The persistent geometry is a coherent rank-two polar flag;
the modular geometry is a Lucas-computed contraction kernel; and the
remaining arithmetic is the Frobenius action on a normalized ordered-root
cover.  Redundancy five is the first exceptional-cover case, while the
later theorems show how that cover is propagated, contained, or
replaced in the levels and carrier strata treated here.

\appendix
\section{Statement-adequacy ledger}\label{app:adequacy}

This appendix reproduces verbatim the formal statements on which the
manuscript places Lean-level trust.  It is intended to make statement
adequacy refereeable without requiring the reader to inspect the full
development.  The source anchors are
\path{RelativeConicArcs/PRSResidualQuadratic.lean} and
\path{RelativeConicArcs/PRSRedundancyNine.lean}; the import and audit
gates are those named in Section~\ref{sec:verification}.

\paragraph{Bottom definitions.}
\begingroup\footnotesize
\begin{verbatim}
def dividedPowerContraction {n : ℕ} (r : R)
    (a : Fin (n + 2) → R) : Fin (n + 1) → R :=
  fun i => a i.succ - r * a i.castSucc

structure ResidualSliceInput (S : Type*) (q : ℕ) where
  isDeep : S → Prop
  exceptional : S → Prop
  integralGenusAtMostOneSlice : S → Prop
  rationalPointOutsideDeletedDivisors :
    q ≥ 53 → S → Prop
  splitSquarefreeWitness : S → Prop
  witnessMakesShallow :
    ∀ {s}, splitSquarefreeWitness s → ¬ isDeep s
  pointGivesWitness :
    ∀ (hq : q ≥ 53) {s}, exceptional s →
      integralGenusAtMostOneSlice s →
      rationalPointOutsideDeletedDivisors hq s →
      splitSquarefreeWitness s

structure PersistentFamilyData (S : Type*) (q : ℕ)
    [Fintype S] [DecidableEq S] where
  tangent : Finset S
  sigma : Finset S
  deep : Finset S
  tangent_sigma_disjoint : Disjoint tangent sigma
  deep_eq : deep = tangent ∪ sigma
  fieldOrderPositive : 0 < q
  tangent_card : tangent.card = q * (q + 1)
  sigma_card_doubled : 2 * sigma.card = q * (q ^ 2 - 1)
  orbitCase : OrbitArithmeticCase
\end{verbatim}
\endgroup

The first definition is the paper's contraction.  The two structures
make geometric integrality, deleted-divisor rational points,
witness production, family identification, and family counts explicit
inputs rather than axioms hidden inside a Boolean certificate.

\paragraph{Headline synthesis terminal.}
\begingroup\footnotesize
\begin{verbatim}
theorem redundancyNineSynthesis {S : Type*} {q : ℕ}
    [Fintype S] [DecidableEq S]
    (hq : q ≥ 53)
    (slice : ResidualSliceInput S q)
    (families : PersistentFamilyData S q)
    (exceptional_iff :
      ∀ s, slice.exceptional s ↔ s ∉ families.deep)
    (persistentDeep :
      ∀ s, s ∈ families.deep → slice.isDeep s)
    (hslice :
      ∀ s, slice.exceptional s →
        slice.integralGenusAtMostOneSlice s)
    (hpoint :
      ∀ s (_hs : slice.exceptional s),
        slice.rationalPointOutsideDeletedDivisors hq s) :
    (∀ s, slice.isDeep s ↔ s ∈ families.deep) ∧
      (∀ s, s ∉ families.deep →
        slice.splitSquarefreeWitness s) ∧
      families.deep.card = q * (q + 1) ^ 2 / 2 ∧
      (projectiveOrbitCount families.orbitCase,
        semilinearOrbitCount families.orbitCase) =
      orbitCountPair families.orbitCase
\end{verbatim}
\endgroup

This terminal is adequate for the final synthesis implication once the paper supplies
the arguments \texttt{exceptional\_iff}, \texttt{persistentDeep},
\texttt{hslice}, and \texttt{hpoint}.  It does not prove those
geometric and coding inputs, nor does it derive the orbit table from a
formal group action.

\paragraph{Axiom audit.}
The guarded audit of fourteen terminals reports no axioms for the
finite witness and orbit-pair terminals; only
\texttt{propext} and \texttt{Quot.sound} for ring normalization; and
only \texttt{propext}, \texttt{Classical.choice}, and
\texttt{Quot.sound} for field/root terminals, the family cardinality,
and the headline synthesis theorem.  It uses no project-specific
axiom, native evaluator, finite-data import, or opaque computational
oracle.

\bibliographystyle{alpha}
\bibliography{refs}

\end{document}
